\documentclass[11pt,oneside]{book}

\usepackage[a4paper,margin=1in]{geometry}
\usepackage{amsmath,amssymb,amsthm}
\usepackage{mathtools}
\usepackage{microtype}
\usepackage{lmodern}
\usepackage[hidelinks]{hyperref}
\usepackage{xcolor}

\definecolor{postechred}{RGB}{166,25,85}
\definecolor{kaistblue}{RGB}{0,65,145}

\usepackage{booktabs}
\usepackage{longtable}
\usepackage{array}

\newcolumntype{L}[1]{>{\raggedright\arraybackslash}p{#1}}

\newtheoremstyle{breakplain}
    {6pt}
    {6pt}
    {\itshape}
    {}
    {\bfseries}
    {}
    {\newline}
    {\thmname{#1}\thmnumber{ #2}\thmnote{. #3}}

\newtheoremstyle{breakdefinition}
    {6pt}
    {6pt}
    {}
    {}
    {\bfseries}
    {}
    {\newline}
    {\thmname{#1}\thmnumber{ #2}\thmnote{. #3}}

\newtheoremstyle{sameline}
    {6pt}
    {6pt}
    {}
    {}
    {\itshape}
    {.}
    {0.5em}
    {\thmname{#1}\thmnumber{ #2}\thmnote{. #3}}

\renewenvironment{proof}[1][Proof]
    {\par\noindent\textit{#1.}\par\nobreak\noindent}
    {\hfill$\square$\par}

\theoremstyle{breakdefinition}
\newtheorem{definition}{Definition}[chapter]

\theoremstyle{breakplain}
\newtheorem{theorem}[definition]{Theorem}

\theoremstyle{breakplain}
\newtheorem{proposition}[definition]{Proposition}
\newtheorem{lemma}[definition]{Lemma}
\newtheorem{corollary}[definition]{Corollary}
\newtheorem{conjecture}[definition]{Conjecture}
\newtheorem{hypothesis}[definition]{Hypothesis}

\theoremstyle{breakdefinition}
\newtheorem{formula}[definition]{Formula}
\newtheorem{fact}[definition]{Fact}
\newtheorem{openproblem}[definition]{Open Problem}

\theoremstyle{sameline}
\newtheorem*{remark}{Remark}
\newtheorem*{idea}{Idea}
\newtheorem*{strategy}{Strategy}
\newtheorem*{tactic}{Tactic}

\newenvironment{problem}[1]
    {\par\vspace{6pt}\noindent\textbf{Problem #1.}\hspace{0.5em}\ignorespaces}
    {\par\vspace{6pt}}

\newenvironment{solution}
    {\par\vspace{6pt}\noindent\textit{Solution.}\hspace{0.5em}\ignorespaces}
    {\par\nobreak\noindent\hfill$\square$\par\vspace{6pt}}

% Standard number systems
\newcommand{\N}{\mathbb{N}}
\newcommand{\Nzero}{\mathbb{N}_0}
\newcommand{\Z}{\mathbb{Z}}
\newcommand{\Q}{\mathbb{Q}}
\newcommand{\R}{\mathbb{R}}
\newcommand{\C}{\mathbb{C}}

% General calculus and analysis notation
\newcommand{\dd}{\,d}
\newcommand{\eps}{\varepsilon}
\DeclareMathOperator{\grad}{grad}
\DeclareMathOperator{\divergence}{div}
\DeclareMathOperator{\curl}{curl}
\DeclareMathOperator{\Jac}{Jac}
\DeclareMathOperator{\area}{area}
\DeclareMathOperator{\vol}{vol}
\DeclareMathOperator{\sgn}{sgn}

% Linear algebra operators
\DeclareMathOperator{\rk}{rk}
\DeclareMathOperator{\tr}{tr}
\DeclareMathOperator{\adj}{adj}
\DeclareMathOperator{\Col}{Col}
\DeclareMathOperator{\Null}{Null}
\DeclareMathOperator{\Row}{Row}
\DeclareMathOperator{\Ker}{Ker}
\DeclareMathOperator{\Image}{Im}
\DeclareMathOperator{\Span}{span}
\DeclareMathOperator{\diag}{diag}

% Number theory and algebra operators
\DeclareMathOperator{\lcm}{lcm}
\DeclareMathOperator{\ord}{ord}
\DeclareMathOperator{\Aut}{Aut}
\DeclareMathOperator{\Gal}{Gal}

% Probability and statistics notation
\newcommand{\Prob}{\mathbb{P}}
\newcommand{\Exp}{\mathbb{E}}
\DeclareMathOperator{\Var}{Var}
\DeclareMathOperator{\Cov}{Cov}
\DeclareMathOperator{\Bin}{Bin}
\DeclareMathOperator{\Poi}{Poi}
\DeclareMathOperator{\Normal}{N}

\begin{document}

    \begin{titlepage}
        \centering

        \vspace*{2.4cm}

        {\Huge\bfseries Story of\\[0.35em]
        How \textcolor{postechred}{POSTECH}\\[0.35em]
        Defeats \textcolor{kaistblue}{KAIST}\par}

        \vspace{0.7cm}

        {\LARGE\bfseries Mathematics for Science Quiz\par}

        \vspace{1.6cm}

        {\large
        Jungwoo Kim\\
        Department of Mathematics\\
        \textcolor{postechred}{Pohang University of Science and Technology}
        \par}

        \vfill

        {\large
        \textcolor{postechred}{POSTECH}--\textcolor{kaistblue}{KAIST} Science War\\[0.7em]
        \textcolor{postechred}{Pohang University of Science and Technology}\\
        vs\\
        \textcolor{kaistblue}{Korea Advanced Institute of Science and Technology}
        \par}

        \vspace{0.3cm}

    \end{titlepage}

    \chapter*{Preface}
    \markboth{Preface}{Preface}
    \addcontentsline{toc}{chapter}{Preface}

    This book is devoted to the Mathematics section of Science Quiz in the
    \textcolor{postechred}{POSTECH}--\textcolor{kaistblue}{KAIST} Science War.
    It collects mathematical concepts, facts to memorize, past problems, and
    expected problems that are useful for competition preparation.

    The main purpose of this text is not merely to record solutions, but to
    build a practical and reliable mathematical toolkit for Science Quiz.
    It includes standard theorems, important terminology, memorable facts about
    mathematicians and mathematical history, and problem-solving strategies that
    may be useful in an actual match.

    I sincerely hope that this book contributes to \textcolor{postechred}{POSTECH}'s victory.
    Good luck, and let the game begin.

    \tableofcontents

    \chapter*{Notation and Conventions}
    \markboth{Notation and Conventions}{Notation and Conventions}
    \addcontentsline{toc}{chapter}{Notation and Conventions}

    Throughout this book, we use standard mathematical notation whenever possible.
    The conventions below are arranged to keep the exposition readable,
    unified, and consistent with the chapter structure of Part I. The guiding
    principle is simple: use official, standard notation first, and introduce
    special notation only when it makes a quiz solution faster or clearer.

    \section*{General Conventions}

    Scalars are written in ordinary italic letters, such as $a,b,c,\lambda,t$.
    Vectors are written in bold lowercase letters, such as
    $\mathbf{u},\mathbf{v},\mathbf{x}$. Matrices are written in bold uppercase
    letters, such as $\mathbf{A},\mathbf{B},\mathbf{P}$. Unless otherwise stated,
    vectors are regarded as column vectors. Named theorems, standard English
    terminology, and conventional symbols are preferred over ad hoc notation.

    \vspace{0.5em}

    {\small
    \renewcommand{\arraystretch}{1.18}
    \begin{longtable}{@{}L{0.24\textwidth}L{0.68\textwidth}@{}}
        \toprule
        Notation & Meaning \\
        \midrule
        \endfirsthead

        \toprule
        Notation & Meaning \\
        \midrule
        \endhead

        \bottomrule
        \endlastfoot

        $\N$ & the set of positive integers, $\{1,2,3,\dots\}$ \\
        $\Nzero$ & the set of nonnegative integers, $\{0,1,2,3,\dots\}$ \\
        $\Z$ & the set of integers \\
        $\Q$ & the set of rational numbers \\
        $\R$ & the set of real numbers \\
        $\C$ & the set of complex numbers \\
        $A\subseteq B$ & $A$ is a subset of $B$ \\
        $A\subsetneq B$ & $A$ is a proper subset of $B$ \\
        $A\setminus B$ & the set difference of $A$ and $B$ \\
        $|A|$ & the cardinality of a finite set $A$ \\
        $\varnothing$ & the empty set \\
        $\lambda$ & usually a scalar, eigenvalue, or parameter, depending on context \\
        $\eps$ & a small positive real number, especially in analysis \\
        $\sim$ & may mean asymptotic equivalence, distribution, or an equivalence relation \\
        $O(\cdot)$ & Big-O notation \\
        $o(\cdot)$ & little-o notation \\
        $\square$ & end of a proof or solution \\
    \end{longtable}
    }

    \section*{Theorem-like Environments}

    Part I uses theorem-like labels according to the mathematical role of each
    item. This is especially useful for Science Quiz preparation: a player should
    be able to tell at a glance whether an item is a definition, a named result,
    a formula to memorize, or an open problem. Statement-type environments such
    as Definition and Theorem put the heading on its own line, while short
    guide-type environments such as Idea, Strategy, Tactic, and Remark continue
    on the same line after the label.

    \vspace{0.5em}

    {\small
    \renewcommand{\arraystretch}{1.18}
    \begin{longtable}{@{}L{0.24\textwidth}L{0.68\textwidth}@{}}
        \toprule
        Label & Usage \\
        \midrule
        \endfirsthead

        \toprule
        Label & Usage \\
        \midrule
        \endhead

        \bottomrule
        \endlastfoot

        Definition & a precise meaning of a term or object \\
        Theorem & a major named result, usually worth memorizing by name \\
        Proposition & a useful standard result that is less central than a theorem \\
        Lemma & a supporting result or technical tool \\
        Corollary & a result that follows quickly from a theorem or proposition \\
        Formula & a computational identity or closed form that should be recognized quickly \\
        Fact & a memorization-oriented mathematical or historical fact \\
        Conjecture & a statement believed to be true but not proved in full generality \\
        Hypothesis & a named assumption or guiding claim, often attached to a major open problem \\
        Open Problem & a famous unsolved problem or a problem whose general form remains open \\
        Idea & a short conceptual explanation of why a result is true or useful \\
        Strategy & a fast quiz-solving method or recognition trick \\
        Tactic & a local, immediately usable trick for a particular type of question \\
        Remark & supplementary context, caution, or memorable information \\
    \end{longtable}
    }

    \section*{Calculus and Vector Calculus}

    Calculus notation includes one-variable calculus, multivariable calculus,
    and vector calculus. The symbol $\nabla$ is read as ``nabla'' or ``del.''

    \vspace{0.5em}

    {\small
    \renewcommand{\arraystretch}{1.18}
    \begin{longtable}{@{}L{0.27\textwidth}L{0.65\textwidth}@{}}
        \toprule
        Notation & Meaning \\
        \midrule
        \endfirsthead

        \toprule
        Notation & Meaning \\
        \midrule
        \endhead

        \bottomrule
        \endlastfoot

        $f'(a)$ & the derivative of $f$ at $a$ \\
        $f^{(n)}(a)$ & the $n$th derivative of $f$ at $a$ \\
        $\displaystyle \frac{\dd y}{\dd x}$ & the derivative of $y$ with respect to $x$ \\
        $\displaystyle \int_a^b f(x)\dd x$ & the definite integral of $f$ from $a$ to $b$ \\
        $C^k(I)$ & the set of $k$ times continuously differentiable functions on $I$ \\
        $C^\infty(I)$ & the set of smooth functions on $I$ \\
        $\nabla f$ or $\grad f$ & the gradient of a scalar-valued function $f$ \\
        $\nabla\cdot\mathbf{F}$ or $\divergence\mathbf{F}$ & the divergence of a vector field $\mathbf{F}$ \\
        $\nabla\times\mathbf{F}$ or $\curl\mathbf{F}$ & the curl of a vector field $\mathbf{F}$ \\
        $\Delta f$ & the Laplacian of $f$, usually $\nabla\cdot\nabla f$ \\
        $D\mathbf{F}$ & the derivative matrix, or Jacobian matrix, of a map $\mathbf{F}$ \\
        $\Jac(\mathbf{F})$ & the Jacobian determinant of $\mathbf{F}$, when defined \\
        $\displaystyle \int_C \mathbf{F}\cdot\dd\mathbf{r}$ & a line integral of a vector field along a curve $C$ \\
        $\displaystyle \iint_S \mathbf{F}\cdot\mathbf{n}\dd S$ & a flux integral across a surface $S$ \\
    \end{longtable}
    }

    \section*{Linear Algebra}

    In linear algebra, we consistently use bold notation for vectors and matrices.
    For example, $\mathbf{v}\in\R^n$ is a vector and
    $\mathbf{A}\in\R^{m\times n}$ is a matrix. The scalar $\lambda$ is an
    eigenvalue, not a vector.

    \vspace{0.5em}

    {\small
    \renewcommand{\arraystretch}{1.18}
    \begin{longtable}{@{}L{0.27\textwidth}L{0.65\textwidth}@{}}
        \toprule
        Notation & Meaning \\
        \midrule
        \endfirsthead

        \toprule
        Notation & Meaning \\
        \midrule
        \endhead

        \bottomrule
        \endlastfoot

        $\mathbf{A}^{\mathsf T}$ & the transpose of $\mathbf{A}$ \\
        $\mathbf{I}_n$ & the $n\times n$ identity matrix \\
        $\mathbf{0}$ & the zero vector or zero matrix, depending on context \\
        $\det(\mathbf{A})$ & the determinant of $\mathbf{A}$ \\
        $\rk(\mathbf{A})$ & the rank of $\mathbf{A}$ \\
        $\tr(\mathbf{A})$ & the trace of $\mathbf{A}$ \\
        $\adj(\mathbf{A})$ & the adjugate of $\mathbf{A}$ \\
        $\Col(\mathbf{A})$ & the column space of $\mathbf{A}$ \\
        $\Null(\mathbf{A})$ & the null space of $\mathbf{A}$ \\
        $\Col(\mathbf{A}^{\mathsf T})$ & the row space of $\mathbf{A}$ \\
        $\Row(\mathbf{A})$ & another notation for the row space, used only when convenient \\
        $\Ker(T)$ & the kernel of a linear map $T$ \\
        $\Image(T)$ & the image of a linear map $T$ \\
        $\Span(\mathbf{v}_1,\dots,\mathbf{v}_k)$ & the span of vectors $\mathbf{v}_1,\dots,\mathbf{v}_k$ \\
        $\langle\mathbf{u},\mathbf{v}\rangle$ & the inner product of $\mathbf{u}$ and $\mathbf{v}$ \\
        $\|\mathbf{v}\|$ & the norm of $\mathbf{v}$ \\
        $\mathbf{u}\perp\mathbf{v}$ & $\mathbf{u}$ and $\mathbf{v}$ are orthogonal \\
        $U\oplus V$ & the direct sum of subspaces $U$ and $V$ \\
    \end{longtable}
    }

    The four fundamental subspaces of $\mathbf{A}\in\R^{m\times n}$ are written as
    \[
        \Col(\mathbf{A}),\qquad
        \Null(\mathbf{A}),\qquad
        \Col(\mathbf{A}^{\mathsf T}),\qquad
        \Null(\mathbf{A}^{\mathsf T}).
    \]
    We use the standard orthogonal decompositions
    \[
        \R^m=\Col(\mathbf{A})\oplus\Null(\mathbf{A}^{\mathsf T}),
        \qquad
        \R^n=\Col(\mathbf{A}^{\mathsf T})\oplus\Null(\mathbf{A}).
    \]

    For cofactors, we write $A_{ij}$ for the scalar cofactor of the entry
    $a_{ij}$ of $\mathbf{A}$. The cofactor matrix is denoted by
    \[
        \mathbf{C}=(A_{ij}),
    \]
    and the adjugate matrix is
    \[
        \adj(\mathbf{A})=\mathbf{C}^{\mathsf T}.
    \]
    We avoid using $\mathbf{A}^{*}$ for the adjugate, since $\mathbf{A}^{*}$
    often denotes the conjugate transpose in other contexts.

    \section*{Differential Equations}

    Differential equations are written using standard derivative notation.
    The independent variable is usually $t$ or $x$, depending on the context.

    \vspace{0.5em}

    {\small
    \renewcommand{\arraystretch}{1.18}
    \begin{longtable}{@{}L{0.24\textwidth}L{0.68\textwidth}@{}}
        \toprule
        Notation & Meaning \\
        \midrule
        \endfirsthead

        \toprule
        Notation & Meaning \\
        \midrule
        \endhead

        \bottomrule
        \endlastfoot

        $y'$ & the first derivative of $y$ with respect to the relevant variable \\
        $y''$ & the second derivative of $y$ with respect to the relevant variable \\
        $\dot{x}$ & the derivative of $x$ with respect to time $t$ \\
        $\ddot{x}$ & the second derivative of $x$ with respect to time $t$ \\
        $y_h$ & the homogeneous part of a solution \\
        $y_p$ & a particular solution \\
        $c_1,c_2,\dots$ & arbitrary constants in a general solution \\
        $\mathcal{L}\{f(t)\}(s)$ & the Laplace transform of $f(t)$ \\
    \end{longtable}
    }

    \section*{Combinatorics, Discrete Mathematics, Probability and Statistics}

    Combinatorics and discrete mathematics provide the notation for counting,
    recurrence relations, and graphs. Probability and statistics use this
    notation frequently, especially in finite probability spaces.

    \vspace{0.5em}

    {\small
    \renewcommand{\arraystretch}{1.18}
    \begin{longtable}{@{}L{0.27\textwidth}L{0.65\textwidth}@{}}
        \toprule
        Notation & Meaning \\
        \midrule
        \endfirsthead

        \toprule
        Notation & Meaning \\
        \midrule
        \endhead

        \bottomrule
        \endlastfoot

        $n!$ & factorial \\
        $\binom{n}{k}$ & the number of $k$-element subsets of an $n$-element set \\
        $P(n,k)$ & the number of $k$-permutations of an $n$-element set \\
        $a_n$ & the $n$th term of a sequence \\
        $G=(V,E)$ & a graph with vertex set $V$ and edge set $E$ \\
        $\deg(v)$ & the degree of a vertex $v$ in a graph \\
        $K_n$ & the complete graph on $n$ vertices \\
        $C_n$ & the cycle graph on $n$ vertices \\
        $P_n$ & the path graph on $n$ vertices \\
        $\Prob(A)$ & the probability of an event $A$ \\
        $\Prob(A\mid B)$ & the conditional probability of $A$ given $B$ \\
        $\Exp[X]$ & the expectation of a random variable $X$ \\
        $\Var(X)$ & the variance of $X$ \\
        $\Cov(X,Y)$ & the covariance of $X$ and $Y$ \\
        $X\sim\Bin(n,p)$ & $X$ follows a binomial distribution with parameters $n$ and $p$ \\
        $X\sim\Poi(\lambda)$ & $X$ follows a Poisson distribution with parameter $\lambda$ \\
        $X\sim\Normal(\mu,\sigma^2)$ & $X$ follows a normal distribution with mean $\mu$ and variance $\sigma^2$ \\
    \end{longtable}
    }

    Random variables are usually denoted by uppercase letters such as $X,Y,Z$.
    Events are denoted by uppercase letters such as $A,B,E$.
    In graph-theoretic contexts, $\deg(v)$ denotes the degree of a vertex.

    \section*{Geometry}

    We use standard geometric notation unless otherwise stated.

    \vspace{0.5em}

    {\small
    \renewcommand{\arraystretch}{1.18}
    \begin{longtable}{@{}L{0.24\textwidth}L{0.68\textwidth}@{}}
        \toprule
        Notation & Meaning \\
        \midrule
        \endfirsthead

        \toprule
        Notation & Meaning \\
        \midrule
        \endhead

        \bottomrule
        \endlastfoot

        $\triangle ABC$ & the triangle with vertices $A,B,C$ \\
        $\angle ABC$ & the angle with vertex $B$ \\
        $\overline{AB}$ & the line segment joining $A$ and $B$ \\
        $AB$ & the length of the segment $\overline{AB}$, when no confusion is possible \\
        $[ABC]$ & the area of triangle $ABC$, when used in geometry problems \\
        $\parallel$ & parallel \\
        $\perp$ & perpendicular \\
    \end{longtable}
    }

    \section*{Number Theory and Algebra}

    Unless otherwise stated, divisibility and congruence statements are made in
    $\Z$. In algebra, the notation is kept standard and lightweight, since
    Science Quiz rarely requires highly specialized algebraic formalism.

    \vspace{0.5em}

    {\small
    \renewcommand{\arraystretch}{1.18}
    \begin{longtable}{@{}L{0.27\textwidth}L{0.65\textwidth}@{}}
        \toprule
        Notation & Meaning \\
        \midrule
        \endfirsthead

        \toprule
        Notation & Meaning \\
        \midrule
        \endhead

        \bottomrule
        \endlastfoot

        $a\mid b$ & $a$ divides $b$ \\
        $a\nmid b$ & $a$ does not divide $b$ \\
        $\gcd(a,b)$ & the greatest common divisor of $a$ and $b$ \\
        $\lcm(a,b)$ & the least common multiple of $a$ and $b$ \\
        $a\equiv b\pmod n$ & $a$ is congruent to $b$ modulo $n$ \\
        $\varphi(n)$ & Euler's phi function, or Euler's totient function \\
        $\ord_n(a)$ & the order of $a$ modulo $n$ \\
        $p$ & usually a prime number, when used in number-theoretic contexts \\
        $G$ & usually a group \\
        $e$ & the identity element of a group, when the context is clear \\
        $H\leq G$ & $H$ is a subgroup of $G$ \\
        $H\trianglelefteq G$ & $H$ is a normal subgroup of $G$ \\
        $G/H$ & the quotient group of $G$ by $H$ \\
        $R$ & usually a ring \\
        $F$ or $K$ & usually a field \\
        $F[x]$ & the polynomial ring over a field $F$ \\
        $\deg f$ & the degree of a polynomial $f$ \\
        $\Aut(G)$ & the automorphism group of $G$ \\
        $\Gal(L/K)$ & the Galois group of a field extension $L/K$ \\
    \end{longtable}
    }

    The Chinese Remainder Theorem may be abbreviated as CRT after it is first
    introduced. In algebraic contexts, $\deg f$ denotes the degree of a polynomial
    $f$; in graph-theoretic contexts, $\deg(v)$ denotes the degree of a vertex.

    \section*{Miscellaneous Topics}

    Some symbols have different meanings in different areas. In this book, the
    intended meaning should always be clear from context.

    \vspace{0.5em}

    {\small
    \renewcommand{\arraystretch}{1.18}
    \begin{longtable}{@{}L{0.27\textwidth}L{0.65\textwidth}@{}}
        \toprule
        Notation & Meaning \\
        \midrule
        \endfirsthead

        \toprule
        Notation & Meaning \\
        \midrule
        \endhead

        \bottomrule
        \endlastfoot

        $\lfloor x\rfloor$ & the floor of $x$ \\
        $\lceil x\rceil$ & the ceiling of $x$ \\
        $\{x\}$ & the fractional part of $x$, usually $x-\lfloor x\rfloor$ \\
        RH & the Riemann Hypothesis, after it has been introduced \\
        P vs NP & the P versus NP problem, after it has been introduced \\
    \end{longtable}
    }

    \part{Mathematical Concepts}

    \chapter{Calculus}

    This chapter collects quiz-friendly results from one-variable calculus,
    multivariable calculus, and vector calculus. The emphasis is on named
    theorems, memorable formulas, and geometric ideas that are easy to ask and
    useful to recognize quickly.

    \section*{One-Variable Calculus}

    \begin{formula}[Standard Trigonometric and Exponential Limits]
        The following limits are fundamental:
        \[
            \lim_{x\to 0}\frac{\sin x}{x}=1,
            \qquad
            \lim_{x\to 0}\frac{1-\cos x}{x^2}=\frac12,
        \]
        \[
            \lim_{x\to 0}\frac{e^x-1}{x}=1,
            \qquad
            \lim_{x\to 0}\frac{\log(1+x)}{x}=1,
        \]
        and
        \[
            \lim_{x\to 0}(1+x)^{1/x}=e.
        \]
    \end{formula}

    \begin{strategy}
        These limits are often faster than L'Hopital's Rule. In a quiz, the
        expressions $\sin x$, $1-\cos x$, $e^x-1$, $\log(1+x)$, and
        $(1+x)^{1/x}$ are strong signals for these standard patterns.
    \end{strategy}

    \begin{theorem}[Intermediate Value Theorem]
        Let $f$ be continuous on $[a,b]$. If $N$ is any number between
        $f(a)$ and $f(b)$, then there exists $c\in[a,b]$ such that
        \[
            f(c)=N.
        \]
    \end{theorem}

    \begin{strategy}
        The Intermediate Value Theorem is the theorem behind many existence
        arguments. To show that an equation has a solution, define a continuous
        function and show that its values cross the desired level.
    \end{strategy}

    \begin{theorem}[Extreme Value Theorem]
        Let $f$ be continuous on $[a,b]$. Then $f$ attains both a maximum and
        a minimum on $[a,b]$.
    \end{theorem}

    \begin{remark}
        The closed interval condition is important. A continuous function on an
        open interval need not attain a maximum or a minimum.
    \end{remark}

    \begin{theorem}[Mean Value Theorem]
        Let $f$ be continuous on $[a,b]$ and differentiable on $(a,b)$. Then
        there exists $c\in(a,b)$ such that
        \[
            f'(c)
            =
            \frac{f(b)-f(a)}{b-a}.
        \]
    \end{theorem}

    \begin{idea}
        The Mean Value Theorem says that somewhere on the interval, the
        instantaneous rate of change equals the average rate of change.
    \end{idea}

    \begin{theorem}[Taylor's Theorem]
        Suppose that $f$ is $n+1$ times differentiable near $a$. Then
        \[
            f(x)
            =
            f(a)+f'(a)(x-a)+\frac{f''(a)}{2!}(x-a)^2
            +\cdots
            +\frac{f^{(n)}(a)}{n!}(x-a)^n
            +R_n(x),
        \]
        where, in Lagrange's form of the remainder,
        \[
            R_n(x)
            =
            \frac{f^{(n+1)}(c)}{(n+1)!}(x-a)^{n+1}
        \]
        for some $c$ between $a$ and $x$.
    \end{theorem}

    \begin{strategy}
        Taylor expansion is one of the fastest tools for limits, approximations,
        and hidden series evaluations. Expressions involving $e^x$, $\sin x$,
        $\cos x$, $\log(1+x)$, or $\arctan x$ often invite Taylor expansion.
    \end{strategy}

    \begin{formula}[Maclaurin Series]
        The most important Maclaurin series are
        \[
            e^x
            =
            \sum_{n=0}^{\infty}\frac{x^n}{n!},
        \]
        \[
            \sin x
            =
            \sum_{n=0}^{\infty}(-1)^n\frac{x^{2n+1}}{(2n+1)!},
            \qquad
            \cos x
            =
            \sum_{n=0}^{\infty}(-1)^n\frac{x^{2n}}{(2n)!},
        \]
        \[
            \frac{1}{1-x}
            =
            \sum_{n=0}^{\infty}x^n
            \qquad (|x|<1),
        \]
        \[
            \log(1+x)
            =
            \sum_{n=1}^{\infty}(-1)^{n+1}\frac{x^n}{n}
            \qquad (-1<x\leq 1),
        \]
        and
        \[
            \arctan x
            =
            \sum_{n=0}^{\infty}(-1)^n\frac{x^{2n+1}}{2n+1}
            \qquad (|x|\leq 1).
        \]
    \end{formula}

    \begin{formula}[Euler's Formula]
        For every real number $x$,
        \[
            e^{ix}
            =
            \cos x+i\sin x.
        \]
        In particular,
        \[
            e^{i\pi}+1=0.
        \]
    \end{formula}

    \begin{remark}
        The identity $e^{i\pi}+1=0$ is called Euler's identity. It connects
        the constants $e$, $i$, $\pi$, $1$, and $0$ in a single formula.
    \end{remark}

    \begin{formula}[Gaussian Integral]
        The Gaussian integral is
        \[
            \int_{-\infty}^{\infty}e^{-x^2}\dd x
            =
            \sqrt{\pi}.
        \]
        Equivalently,
        \[
            \int_0^\infty e^{-x^2}\dd x
            =
            \frac{\sqrt{\pi}}{2}.
        \]
    \end{formula}

    \begin{remark}
        The Gaussian integral is one of the most famous integrals in calculus.
        It is closely connected to the normal distribution in probability.
    \end{remark}

    \begin{formula}[Wallis Product]
        The Wallis product is
        \[
            \frac{\pi}{2}
            =
            \prod_{n=1}^{\infty}
            \frac{(2n)(2n)}{(2n-1)(2n+1)}
            =
            \frac{2}{1}\cdot\frac{2}{3}\cdot
            \frac{4}{3}\cdot\frac{4}{5}\cdot
            \frac{6}{5}\cdot\frac{6}{7}\cdots.
        \]
    \end{formula}

    \begin{definition}[Gamma Function]
        For $x>0$, the Gamma function is defined by
        \[
            \Gamma(x)
            =
            \int_0^\infty t^{x-1}e^{-t}\dd t.
        \]
    \end{definition}

    \begin{proposition}[Gamma Function and Factorials]
        The Gamma function satisfies
        \[
            \Gamma(x+1)=x\Gamma(x).
        \]
        In particular, for every positive integer $n$,
        \[
            \Gamma(n)=(n-1)!.
        \]
        Also,
        \[
            \Gamma\left(\frac12\right)=\sqrt{\pi}.
        \]
    \end{proposition}

    \begin{formula}[Stirling's Formula]
        As $n\to\infty$,
        \[
            n!
            \sim
            \sqrt{2\pi n}\left(\frac{n}{e}\right)^n.
        \]
    \end{formula}

    \begin{remark}
        Stirling's formula is the standard approximation for large factorials.
        It often appears in counting, probability, and asymptotic estimation.
    \end{remark}

    \section*{Multivariable Calculus and Vector Calculus}

    \begin{definition}[Gradient, Divergence, Curl, and Laplacian]
        For a scalar-valued function $f(x,y,z)$ and a vector field
        $\mathbf{F}=(P,Q,R)$, we write
        \[
            \nabla f
            =
            \left(\frac{\partial f}{\partial x},
                  \frac{\partial f}{\partial y},
                  \frac{\partial f}{\partial z}\right),
        \]
        \[
            \nabla\cdot\mathbf{F}
            =
            \frac{\partial P}{\partial x}
            +
            \frac{\partial Q}{\partial y}
            +
            \frac{\partial R}{\partial z},
        \]
        \[
            \nabla\times\mathbf{F}
            =
            \begin{vmatrix}
                \mathbf{i} & \mathbf{j} & \mathbf{k} \\
                \partial/\partial x & \partial/\partial y & \partial/\partial z \\
                P & Q & R
            \end{vmatrix},
        \]
        and
        \[
            \Delta f
            =
            \nabla\cdot\nabla f
            =
            f_{xx}+f_{yy}+f_{zz}.
        \]
    \end{definition}

    \begin{idea}
        The gradient points in the direction of steepest increase. Divergence
        measures source or sink behavior. Curl measures local rotation. The
        Laplacian measures how a value differs from its surrounding average.
    \end{idea}

    \begin{theorem}[Clairaut's Theorem]
        If $f$ has continuous second partial derivatives near a point, then the
        mixed partial derivatives agree at that point:
        \[
            \frac{\partial^2 f}{\partial x\partial y}
            =
            \frac{\partial^2 f}{\partial y\partial x}.
        \]
    \end{theorem}

    \begin{definition}[Jacobian]
        For a map
        \[
            \mathbf{F}(u_1,\dots,u_n)
            =
            (x_1(u_1,\dots,u_n),\dots,x_n(u_1,\dots,u_n)),
        \]
        the Jacobian matrix is
        \[
            D\mathbf{F}
            =
            \left(\frac{\partial x_i}{\partial u_j}\right)_{i,j=1}^{n}.
        \]
        Its determinant
        \[
            \Jac(\mathbf{F})
            =
            \det(D\mathbf{F})
        \]
        is called the Jacobian determinant.
    \end{definition}

    \begin{theorem}[Change of Variables Formula]
        Under suitable regularity and one-to-one assumptions,
        \[
            \int_{\mathbf{F}(D)} f(\mathbf{x})\dd\mathbf{x}
            =
            \int_D f(\mathbf{F}(\mathbf{u}))
            \left|\det D\mathbf{F}(\mathbf{u})\right|\dd\mathbf{u}.
        \]
    \end{theorem}

    \begin{remark}
        The Jacobian determinant is the local area or volume scaling factor of
        a change of variables.
    \end{remark}

    \begin{theorem}[Lagrange Multipliers]
        Suppose that $f$ has a local maximum or minimum subject to the constraint
        \[
            g(x_1,\dots,x_n)=0
        \]
        at a point where $\nabla g\neq\mathbf{0}$. Then there exists a scalar
        $\lambda$ such that
        \[
            \nabla f=\lambda\nabla g.
        \]
    \end{theorem}

    \begin{idea}
        At a constrained optimum, the level surface of $f$ is tangent to the
        constraint surface. Thus their normal vectors are parallel.
    \end{idea}

    \begin{theorem}[Green's Theorem]
        Let $C$ be a positively oriented, simple closed curve in the plane, and
        let $D$ be the region enclosed by $C$. Then
        \[
            \oint_C P\dd x+Q\dd y
            =
            \iint_D
            \left(
                \frac{\partial Q}{\partial x}
                -
                \frac{\partial P}{\partial y}
            \right)
            \dd A.
        \]
    \end{theorem}

    \begin{idea}
        Green's Theorem converts circulation along a boundary curve into curl
        over the region inside. It is the two-dimensional ancestor of Stokes'
        Theorem.
    \end{idea}

    \begin{theorem}[Stokes' Theorem]
        Let $S$ be an oriented smooth surface with positively oriented boundary
        curve $\partial S$. Then
        \[
            \oint_{\partial S}\mathbf{F}\cdot\dd\mathbf{r}
            =
            \iint_S (\nabla\times\mathbf{F})\cdot\mathbf{n}\dd S.
        \]
    \end{theorem}

    \begin{idea}
        Stokes' Theorem says that circulation around the boundary equals total
        curl through the surface. It turns a boundary integral into an interior
        integral.
    \end{idea}

    \begin{theorem}[Divergence Theorem]
        Let $E$ be a solid region in $\R^3$ with outward-oriented boundary
        surface $\partial E$. Then
        \[
            \iint_{\partial E}\mathbf{F}\cdot\mathbf{n}\dd S
            =
            \iiint_E \nabla\cdot\mathbf{F}\dd V.
        \]
    \end{theorem}

    \begin{idea}
        The Divergence Theorem says that outward flux through the boundary
        equals total source strength inside the region.
    \end{idea}

    \begin{remark}
        Green's Theorem, Stokes' Theorem, and the Divergence Theorem share the
        same philosophy: a boundary integral is equal to an interior integral.
        This is the geometric heart of vector calculus.
    \end{remark}

    \chapter{Linear Algebra}

    This chapter focuses on named theorems, matrix invariants, and geometric
    principles that are especially likely to appear in quiz-style questions.

    \begin{theorem}[Rank--Nullity Theorem]
        Let $T:V\to W$ be a linear map between finite-dimensional vector spaces.
        Then
        \[
            \dim V
            =
            \dim\Ker(T)+\dim\Image(T).
        \]
        For a matrix $\mathbf{A}\in F^{m\times n}$, this becomes
        \[
            n=\rk(\mathbf{A})+\dim\Null(\mathbf{A}).
        \]
    \end{theorem}

    \begin{strategy}
        Rank--Nullity connects variables, pivots, and free variables. It is the
        linear algebra version of accounting: nothing disappears; dimensions
        are redistributed.
    \end{strategy}

    \begin{theorem}[Invertible Matrix Theorem]
        For a square matrix $\mathbf{A}\in F^{n\times n}$, the following
        statements are equivalent:
        \[
            \mathbf{A}\text{ is invertible},
            \qquad
            \det(\mathbf{A})\neq 0,
            \qquad
            \rk(\mathbf{A})=n,
        \]
        \[
            \Null(\mathbf{A})=\{\mathbf{0}\},
        \]
        and
        \[
            \mathbf{A}\mathbf{x}=\mathbf{b}
        \]
        has a unique solution for every $\mathbf{b}\in F^n$.
    \end{theorem}

    \begin{remark}
        This theorem is a dictionary. It translates between determinants, rank,
        null spaces, and solvability of linear systems.
    \end{remark}

    \begin{definition}[Determinant]
        The determinant is a scalar associated to a square matrix
        $\mathbf{A}\in F^{n\times n}$, denoted by $\det(\mathbf{A})$.
        Geometrically, for real matrices, $|\det(\mathbf{A})|$ is the volume
        scaling factor of the linear transformation
        $\mathbf{x}\mapsto\mathbf{A}\mathbf{x}$.
    \end{definition}

    \begin{proposition}[Basic Properties of Determinants]
        For square matrices $\mathbf{A},\mathbf{B}\in F^{n\times n}$, we have
        \[
            \det(\mathbf{A}\mathbf{B})
            =
            \det(\mathbf{A})\det(\mathbf{B}),
        \]
        \[
            \det(\mathbf{A}^{\mathsf T})=\det(\mathbf{A}),
        \]
        and, if $\mathbf{A}$ is invertible,
        \[
            \det(\mathbf{A}^{-1})
            =
            \frac{1}{\det(\mathbf{A})}.
        \]
    \end{proposition}

    \begin{formula}[Determinant of a Triangular Matrix]
        If $\mathbf{A}$ is upper triangular or lower triangular with diagonal
        entries $a_{11},a_{22},\dots,a_{nn}$, then
        \[
            \det(\mathbf{A})
            =
            a_{11}a_{22}\cdots a_{nn}.
        \]
    \end{formula}

    \begin{proposition}[Trace, Determinant, and Eigenvalues]
        Let $\mathbf{A}\in\C^{n\times n}$ have eigenvalues
        $\lambda_1,\dots,\lambda_n$, counted with algebraic multiplicity. Then
        \[
            \tr(\mathbf{A})
            =
            \lambda_1+\cdots+\lambda_n
        \]
        and
        \[
            \det(\mathbf{A})
            =
            \lambda_1\lambda_2\cdots\lambda_n.
        \]
    \end{proposition}

    \begin{definition}[Eigenvalue and Characteristic Polynomial]
        Let $\mathbf{A}\in F^{n\times n}$. A scalar $\lambda\in F$ is an
        eigenvalue of $\mathbf{A}$ if there exists a nonzero vector
        $\mathbf{v}\in F^n$ such that
        \[
            \mathbf{A}\mathbf{v}=\lambda\mathbf{v}.
        \]
        The characteristic polynomial of $\mathbf{A}$ is
        \[
            p_{\mathbf{A}}(\lambda)
            =
            \det(\lambda\mathbf{I}_n-\mathbf{A}).
        \]
    \end{definition}

    \begin{proposition}[Eigenvalues and the Characteristic Polynomial]
        A scalar $\lambda$ is an eigenvalue of $\mathbf{A}$ if and only if
        \[
            \det(\lambda\mathbf{I}_n-\mathbf{A})=0.
        \]
        Equivalently, the eigenvalues of $\mathbf{A}$ are the roots of the
        characteristic polynomial $p_{\mathbf{A}}(\lambda)$.
    \end{proposition}

    \begin{theorem}[Cayley--Hamilton Theorem]
        Every square matrix satisfies its own characteristic polynomial. That
        is, if
        \[
            p_{\mathbf{A}}(\lambda)
            =
            \det(\lambda\mathbf{I}_n-\mathbf{A}),
        \]
        then
        \[
            p_{\mathbf{A}}(\mathbf{A})=\mathbf{0}.
        \]
    \end{theorem}

    \begin{remark}
        Cayley--Hamilton is a famous named theorem. It is especially useful for
        reducing high powers of a matrix to lower powers.
    \end{remark}

    \begin{theorem}[Diagonalization Criterion]
        A matrix $\mathbf{A}\in F^{n\times n}$ is diagonalizable over $F$ if
        and only if $F^n$ has a basis consisting of eigenvectors of
        $\mathbf{A}$.
    \end{theorem}

    \begin{corollary}[Distinct Eigenvalues Imply Diagonalizability]
        If an $n\times n$ matrix $\mathbf{A}$ has $n$ distinct eigenvalues in
        $F$, then $\mathbf{A}$ is diagonalizable over $F$.
    \end{corollary}

    \begin{strategy}
        Distinct eigenvalues are an easy sufficient condition for
        diagonalizability. They are not necessary: a matrix can be
        diagonalizable even when some eigenvalues are repeated.
    \end{strategy}

    \begin{theorem}[Spectral Theorem for Real Symmetric Matrices]
        If $\mathbf{A}\in\R^{n\times n}$ is symmetric, then all eigenvalues of
        $\mathbf{A}$ are real, and there exists an orthogonal matrix
        $\mathbf{Q}$ and a diagonal matrix $\mathbf{D}$ such that
        \[
            \mathbf{A}
            =
            \mathbf{Q}\mathbf{D}\mathbf{Q}^{\mathsf T}.
        \]
        Equivalently, $\R^n$ has an orthonormal basis consisting of eigenvectors
        of $\mathbf{A}$.
    \end{theorem}

    \begin{remark}
        In quiz settings, the words ``real symmetric matrix'' are a strong
        signal for real eigenvalues and orthogonal diagonalization.
    \end{remark}

    \begin{theorem}[Cauchy--Schwarz Inequality]
        For vectors $\mathbf{u},\mathbf{v}\in\R^n$, we have
        \[
            |\langle\mathbf{u},\mathbf{v}\rangle|
            \leq
            \|\mathbf{u}\|\,\|\mathbf{v}\|.
        \]
        Equivalently, for real numbers $a_1,\dots,a_n$ and
        $b_1,\dots,b_n$,
        \[
            \left(\sum_{i=1}^{n}a_i b_i\right)^2
            \leq
            \left(\sum_{i=1}^{n}a_i^2\right)
            \left(\sum_{i=1}^{n}b_i^2\right).
        \]
    \end{theorem}

    \begin{strategy}
        The expression $\sum a_i b_i$ is a strong signal for Cauchy--Schwarz.
        The vector form is often the fastest way to recognize the inequality.
    \end{strategy}

    \begin{corollary}[Triangle Inequality]
        For vectors $\mathbf{u},\mathbf{v}\in\R^n$, we have
        \[
            \|\mathbf{u}+\mathbf{v}\|
            \leq
            \|\mathbf{u}\|+\|\mathbf{v}\|.
        \]
    \end{corollary}

    \begin{theorem}[Gram--Schmidt Process]
        From any linearly independent list
        \[
            \mathbf{v}_1,\dots,\mathbf{v}_k
        \]
        in an inner product space, one can construct an orthonormal list
        \[
            \mathbf{q}_1,\dots,\mathbf{q}_k
        \]
        with the same span:
        \[
            \Span(\mathbf{q}_1,\dots,\mathbf{q}_k)
            =
            \Span(\mathbf{v}_1,\dots,\mathbf{v}_k).
        \]
    \end{theorem}

    \begin{remark}
        Gram--Schmidt is the standard procedure for turning an arbitrary basis
        into an orthonormal basis.
    \end{remark}

    \begin{formula}[Projection Matrix]
        If the columns of $\mathbf{Q}\in\R^{n\times k}$ are orthonormal, then
        the orthogonal projection matrix onto $\Col(\mathbf{Q})$ is
        \[
            \mathbf{P}
            =
            \mathbf{Q}\mathbf{Q}^{\mathsf T}.
        \]
        More generally, if the columns of $\mathbf{A}\in\R^{n\times k}$ are
        linearly independent, then the orthogonal projection matrix onto
        $\Col(\mathbf{A})$ is
        \[
            \mathbf{P}
            =
            \mathbf{A}(\mathbf{A}^{\mathsf T}\mathbf{A})^{-1}
            \mathbf{A}^{\mathsf T}.
        \]
    \end{formula}

    \begin{proposition}[Projection Matrix Criterion]
        A matrix $\mathbf{P}$ is a projection matrix if
        \[
            \mathbf{P}^2=\mathbf{P}.
        \]
        It is an orthogonal projection matrix if, in addition,
        \[
            \mathbf{P}^{\mathsf T}=\mathbf{P}.
        \]
    \end{proposition}

    \begin{definition}[Singular Values]
        Let $\mathbf{A}\in\R^{m\times n}$. The singular values of
        $\mathbf{A}$ are the nonnegative square roots of the eigenvalues of
        \[
            \mathbf{A}^{\mathsf T}\mathbf{A}.
        \]
    \end{definition}

    \begin{theorem}[Singular Value Decomposition]
        Every real matrix $\mathbf{A}\in\R^{m\times n}$ can be written as
        \[
            \mathbf{A}
            =
            \mathbf{U}\boldsymbol{\Sigma}\mathbf{V}^{\mathsf T},
        \]
        where $\mathbf{U}$ and $\mathbf{V}$ are orthogonal matrices and
        $\boldsymbol{\Sigma}$ is a rectangular diagonal matrix whose diagonal
        entries are the singular values of $\mathbf{A}$.
    \end{theorem}

    \begin{remark}
        Singular Value Decomposition, or SVD, is one of the most powerful
        decompositions in linear algebra. It applies to every real matrix, not
        only to square or diagonalizable matrices.
    \end{remark}

    \begin{theorem}[Jordan Canonical Form]
        Over the complex numbers, every square matrix is similar to a Jordan
        matrix. That is, for every $\mathbf{A}\in\C^{n\times n}$, there exists
        an invertible matrix $\mathbf{P}$ such that
        \[
            \mathbf{P}^{-1}\mathbf{A}\mathbf{P}
        \]
        is in Jordan canonical form.
    \end{theorem}

    \begin{remark}
        Jordan canonical form is usually more important as a named concept than
        as a computational tool in Science Quiz. It describes what remains when
        a matrix cannot be fully diagonalized.
    \end{remark}

    \chapter{Differential Equations}

    This chapter collects named principles and famous model equations from
    ordinary and partial differential equations.

    \begin{theorem}[Picard--Lindelof Existence and Uniqueness Theorem]
        Consider the initial value problem
        \[
            y'=f(t,y),
            \qquad
            y(t_0)=y_0.
        \]
        If $f$ and $\partial f/\partial y$ are continuous near
        $(t_0,y_0)$, then there exists a unique solution near $t=t_0$.
    \end{theorem}

    \begin{remark}
        In a quiz setting, the important phrase is ``existence and uniqueness.''
        It tells us when an initial condition determines exactly one solution.
    \end{remark}

    \begin{theorem}[Superposition Principle]
        If $y_1,\dots,y_k$ are solutions of a homogeneous linear differential
        equation, then any linear combination
        \[
            c_1y_1+\cdots+c_ky_k
        \]
        is also a solution.
    \end{theorem}

    \begin{remark}
        The Superposition Principle is the differential-equation version of
        linear algebra. It is one of the main reasons linear equations are much
        easier to handle than nonlinear equations.
    \end{remark}

    \begin{proposition}[General Solution of a Nonhomogeneous Linear Equation]
        For a nonhomogeneous linear differential equation, the general solution
        has the form
        \[
            y=y_h+y_p,
        \]
        where $y_h$ is the general solution of the associated homogeneous
        equation and $y_p$ is one particular solution of the nonhomogeneous
        equation.
    \end{proposition}

    \begin{formula}[Second-Order Constant-Coefficient Equation]
        Consider the homogeneous equation
        \[
            ay''+by'+cy=0,
            \qquad
            a\neq 0.
        \]
        Its characteristic equation is
        \[
            ar^2+br+c=0.
        \]
        If the roots are distinct real numbers $r_1,r_2$, then
        \[
            y_h=c_1e^{r_1t}+c_2e^{r_2t}.
        \]
        If the equation has a repeated real root $r$, then
        \[
            y_h=(c_1+c_2t)e^{rt}.
        \]
        If the roots are $\alpha\pm i\beta$ with $\beta\neq 0$, then
        \[
            y_h=e^{\alpha t}
            \left(
                c_1\cos \beta t+c_2\sin \beta t
            \right).
        \]
    \end{formula}

    \begin{remark}
        The characteristic equation converts a constant-coefficient linear ODE
        into an algebraic equation. This is the key trick.
    \end{remark}

    \begin{definition}[Wronskian]
        For two differentiable functions $y_1$ and $y_2$, their Wronskian is
        \[
            W(y_1,y_2)(t)
            =
            \begin{vmatrix}
                y_1(t) & y_2(t) \\
                y_1'(t) & y_2'(t)
            \end{vmatrix}
            =
            y_1(t)y_2'(t)-y_1'(t)y_2(t).
        \]
    \end{definition}

    \begin{proposition}[Wronskian and Linear Independence]
        If
        \[
            W(y_1,y_2)(t_0)\neq 0
        \]
        for some $t_0$, then $y_1$ and $y_2$ are linearly independent.
    \end{proposition}

    \begin{formula}[Exponential Growth and Decay]
        The equation
        \[
            y'=ky
        \]
        has solution
        \[
            y(t)=Ce^{kt}.
        \]
        If $k>0$, this models exponential growth. If $k<0$, this models
        exponential decay.
    \end{formula}

    \begin{formula}[Logistic Equation]
        The logistic equation is
        \[
            y'=ky\left(1-\frac{y}{K}\right),
        \]
        where $k>0$ is the growth rate and $K>0$ is the carrying capacity.
        Its equilibrium solutions are
        \[
            y=0
            \qquad\text{and}\qquad
            y=K.
        \]
    \end{formula}

    \begin{remark}
        The logistic equation is the standard model for population growth with
        limited resources. The parameter $K$ represents the limiting population.
    \end{remark}

    \begin{formula}[Simple Harmonic Oscillator]
        The equation
        \[
            x''+\omega^2x=0
        \]
        has general solution
        \[
            x(t)=c_1\cos \omega t+c_2\sin \omega t.
        \]
        Equivalently,
        \[
            x(t)=A\cos(\omega t-\phi)
        \]
        for suitable constants $A$ and $\phi$.
    \end{formula}

    \begin{remark}
        This is the equation behind ideal mass--spring motion and many small
        oscillation models in physics.
    \end{remark}

    \begin{formula}[Damped Harmonic Oscillator]
        The equation
        \[
            mx''+cx'+kx=0
        \]
        models a damped harmonic oscillator. Its characteristic equation is
        \[
            mr^2+cr+k=0.
        \]
        The motion is underdamped, critically damped, or overdamped according
        as
        \[
            c^2-4mk<0,
            \qquad
            c^2-4mk=0,
            \qquad
            c^2-4mk>0.
        \]
    \end{formula}

    \begin{definition}[Laplace Transform]
        The Laplace transform of a function $f(t)$ is
        \[
            \mathcal{L}\{f(t)\}(s)
            =
            \int_0^\infty e^{-st}f(t)\dd t.
        \]
    \end{definition}

    \begin{proposition}[Basic Laplace Transforms]
        The most common Laplace transforms are
        \[
            \mathcal{L}\{1\}(s)=\frac{1}{s},
            \qquad
            \mathcal{L}\{e^{at}\}(s)=\frac{1}{s-a},
        \]
        \[
            \mathcal{L}\{\sin at\}(s)=\frac{a}{s^2+a^2},
            \qquad
            \mathcal{L}\{\cos at\}(s)=\frac{s}{s^2+a^2}.
        \]
    \end{proposition}

    \begin{remark}
        The Laplace transform turns many differential equations into algebraic
        equations. This is its main computational advantage.
    \end{remark}

    \begin{definition}[Heat Equation]
        The heat equation is the partial differential equation
        \[
            u_t=\alpha u_{xx},
            \qquad
            \alpha>0.
        \]
        It models diffusion, especially the flow of heat in a one-dimensional
        medium.
    \end{definition}

    \begin{definition}[Wave Equation]
        The one-dimensional wave equation is
        \[
            u_{tt}=c^2u_{xx}.
        \]
        It models waves traveling with speed $c$.
    \end{definition}

    \begin{theorem}[d'Alembert's Formula]
        The solution of the wave equation
        \[
            u_{tt}=c^2u_{xx}
        \]
        on the real line with initial conditions
        \[
            u(x,0)=f(x),
            \qquad
            u_t(x,0)=g(x),
        \]
        is
        \[
            u(x,t)
            =
            \frac{f(x-ct)+f(x+ct)}{2}
            +
            \frac{1}{2c}
            \int_{x-ct}^{x+ct}g(s)\dd s.
        \]
    \end{theorem}

    \begin{remark}
        d'Alembert's formula shows that waves move along the characteristic
        lines $x-ct=\text{constant}$ and $x+ct=\text{constant}$.
    \end{remark}

    \begin{definition}[Laplace Equation]
        Laplace's equation is
        \[
            \Delta u=0.
        \]
        In two variables, this means
        \[
            u_{xx}+u_{yy}=0.
        \]
        A function satisfying Laplace's equation is called harmonic.
    \end{definition}

    \begin{definition}[Poisson Equation]
        Poisson's equation is
        \[
            \Delta u=f.
        \]
        It is an inhomogeneous version of Laplace's equation.
    \end{definition}

    \begin{remark}
        The heat equation, wave equation, and Laplace equation are three of the
        most famous PDEs. A quiz problem may ask for their names, forms, or
        physical interpretations.
    \end{remark}

    \chapter{Combinatorics and Discrete Mathematics}

    This chapter collects high-yield results from combinatorics and discrete
    mathematics. The emphasis is on named principles, memorable formulas, and
    results with clear quiz-style answers.

    \begin{formula}[Binomial Theorem; Newton's Binomial Formula]
        For every nonnegative integer $n$,
        \[
            (x+y)^n
            =
            \sum_{k=0}^{n}\binom{n}{k}x^{n-k}y^k.
        \]
        In particular,
        \[
            (1+x)^n
            =
            \sum_{k=0}^{n}\binom{n}{k}x^k.
        \]
    \end{formula}

    \begin{remark}
        The coefficients $\binom{n}{k}$ are the entries of Pascal's triangle.
        This is one of the most recognizable bridges between algebra and
        counting.
    \end{remark}

    \begin{theorem}[Pigeonhole Principle; Dirichlet's Box Principle]
        If $n+1$ objects are placed into $n$ boxes, then at least one box
        contains at least two objects.
        More generally, if $N$ objects are placed into $k$ boxes, then at least
        one box contains at least
        \[
            \left\lceil\frac{N}{k}\right\rceil
        \]
        objects.
    \end{theorem}

    \begin{strategy}
        The Pigeonhole Principle is often hidden in problems asking us to prove
        that two objects must share a property. The hard part is usually
        choosing the right ``boxes.''
    \end{strategy}

    \begin{theorem}[Inclusion--Exclusion Principle; Sieve Formula]
        For finite sets $A_1,\dots,A_n$,
        \[
            \left|\bigcup_{i=1}^{n}A_i\right|
            =
            \sum_i |A_i|
            -
            \sum_{i<j}|A_i\cap A_j|
            +
            \sum_{i<j<k}|A_i\cap A_j\cap A_k|
            -\cdots
            +(-1)^{n+1}|A_1\cap\cdots\cap A_n|.
        \]
    \end{theorem}

    \begin{remark}
        Inclusion--Exclusion is the standard correction mechanism for
        overcounting. It is also historically connected to sieve methods.
    \end{remark}

    \begin{formula}[Derangement Formula]
        Let $D_n$ be the number of derangements of $n$ objects, that is,
        permutations with no fixed points. Then
        \[
            D_n
            =
            n!\sum_{k=0}^{n}\frac{(-1)^k}{k!}.
        \]
        Moreover,
        \[
            D_n
            =
            \left\lfloor \frac{n!}{e}+\frac12\right\rfloor.
        \]
    \end{formula}

    \begin{remark}
        The approximation $D_n\approx n!/e$ is very memorable. In probability
        language, a random permutation has no fixed point with probability
        approximately $1/e$.
    \end{remark}

    \begin{formula}[Catalan Numbers]
        The $n$th Catalan number is
        \[
            C_n
            =
            \frac{1}{n+1}\binom{2n}{n}.
        \]
        It begins as
        \[
            1,\ 1,\ 2,\ 5,\ 14,\ 42,\ 132,\dots.
        \]
    \end{formula}

    \begin{remark}
        Catalan numbers count many different objects, including correct
        parenthesizations, Dyck paths, triangulations of a polygon, and rooted
        full binary trees. The number $42$ is a classic giveaway.
    \end{remark}

    \begin{theorem}[Cayley's Formula]
        The number of labeled trees on $n$ vertices is
        \[
            n^{n-2}.
        \]
    \end{theorem}

    \begin{remark}
        Cayley's formula is a perfect Science Quiz result: the statement is
        short, the answer is surprising, and the name is famous.
    \end{remark}

    \begin{theorem}[Burnside's Lemma; Cauchy--Frobenius Lemma]
        Suppose a finite group $G$ acts on a finite set $X$. The number of
        orbits is
        \[
            \frac{1}{|G|}
            \sum_{g\in G}|\operatorname{Fix}(g)|,
        \]
        where
        \[
            \operatorname{Fix}(g)
            =
            \{x\in X:g\cdot x=x\}.
        \]
    \end{theorem}

    \begin{idea}
        Burnside's Lemma says that the number of essentially different objects
        is the average number of objects fixed by a symmetry.
    \end{idea}

    \begin{lemma}[Handshaking Lemma]
        For a finite undirected graph $G=(V,E)$,
        \[
            \sum_{v\in V}\deg(v)=2|E|.
        \]
        In particular, the number of vertices of odd degree is even.
    \end{lemma}

    \begin{remark}
        The name comes from the idea that every handshake contributes $2$ to
        the total count of hands shaken.
    \end{remark}

    \begin{theorem}[Euler's Formula for Planar Graphs]
        If a connected planar graph has $V$ vertices, $E$ edges, and $F$ faces,
        then
        \[
            V-E+F=2.
        \]
    \end{theorem}

    \begin{remark}
        This is the graph-theoretic version of Euler's polyhedron formula.
        It is one of the most famous formulas involving planar graphs.
    \end{remark}

    \begin{theorem}[Eulerian Trail Criterion]
        A connected finite graph has an Eulerian circuit if and only if every
        vertex has even degree. It has an Eulerian trail but not an Eulerian
        circuit if and only if exactly two vertices have odd degree.
    \end{theorem}

    \begin{remark}
        This theorem is historically connected to Euler's solution of the
        Seven Bridges of Konigsberg problem.
    \end{remark}

    \begin{theorem}[Hall's Marriage Theorem]
        Let $G=(X\cup Y,E)$ be a bipartite graph. There is a matching that
        covers every vertex of $X$ if and only if, for every subset
        $S\subseteq X$,
        \[
            |N(S)|\geq |S|,
        \]
        where $N(S)$ is the set of neighbors of vertices in $S$.
    \end{theorem}

    \begin{strategy}
        Hall's Marriage Theorem is the standard named theorem for proving that
        a perfect assignment or matching exists.
    \end{strategy}

    \begin{theorem}[Ramsey's Theorem]
        For any positive integers $r$ and $s$, there exists an integer
        $R(r,s)$ such that every red-blue coloring of the edges of a sufficiently
        large complete graph contains either a red $K_r$ or a blue $K_s$.
    \end{theorem}

    \begin{formula}[The Party Problem]
        The classical Ramsey number
        \[
            R(3,3)=6
        \]
        says that among any six people, there are either three mutual
        acquaintances or three mutual strangers.
    \end{formula}

    \begin{remark}
        Ramsey theory is the mathematics of unavoidable structure. Its slogan
        is that complete disorder is impossible in a large enough system.
    \end{remark}

    \begin{theorem}[Erdos--Szekeres Theorem]
        Any sequence of
        \[
            (r-1)(s-1)+1
        \]
        distinct real numbers contains either an increasing subsequence of
        length $r$ or a decreasing subsequence of length $s$.
    \end{theorem}

    \begin{remark}
        This is a beautiful pigeonhole-style result. It is often remembered as
        the monotone subsequence theorem.
    \end{remark}

    \chapter{Probability and Statistics}

    This chapter focuses on probability results that are easy to ask in a quiz:
    named theorems, famous distributions, memorable approximations, and
    inequalities with clear meanings.

    \begin{theorem}[Law of Total Probability]
        If $B_1,\dots,B_n$ form a partition of the sample space and
        $\Prob(B_i)>0$ for all $i$, then
        \[
            \Prob(A)
            =
            \sum_{i=1}^{n}\Prob(A\mid B_i)\Prob(B_i).
        \]
    \end{theorem}

    \begin{theorem}[Bayes' Theorem]
        If $B_1,\dots,B_n$ form a partition of the sample space and
        $\Prob(A)>0$, then
        \[
            \Prob(B_j\mid A)
            =
            \frac{\Prob(A\mid B_j)\Prob(B_j)}
                 {\sum_{i=1}^{n}\Prob(A\mid B_i)\Prob(B_i)}.
        \]
        In the two-event form,
        \[
            \Prob(B\mid A)
            =
            \frac{\Prob(A\mid B)\Prob(B)}{\Prob(A)}.
        \]
    \end{theorem}

    \begin{remark}
        Bayes' Theorem is the standard formula for reversing conditional
        probability. It updates a prior probability after observing evidence.
    \end{remark}

    \begin{proposition}[Linearity of Expectation]
        For random variables $X_1,\dots,X_n$,
        \[
            \Exp[X_1+\cdots+X_n]
            =
            \Exp[X_1]+\cdots+\Exp[X_n].
        \]
        This holds even when the random variables are not independent.
    \end{proposition}

    \begin{strategy}
        Linearity of expectation is often the fastest way to count an expected
        number of objects. Introduce indicator random variables and add their
        expectations.
    \end{strategy}

    \begin{formula}[Variance and Covariance Identities]
        For a random variable $X$,
        \[
            \Var(X)
            =
            \Exp[X^2]-(\Exp[X])^2.
        \]
        For random variables $X$ and $Y$,
        \[
            \Var(X+Y)
            =
            \Var(X)+\Var(Y)+2\Cov(X,Y).
        \]
        If $X$ and $Y$ are independent, then
        \[
            \Var(X+Y)=\Var(X)+\Var(Y).
        \]
    \end{formula}

    \begin{formula}[Binomial Distribution]
        If $X\sim\Bin(n,p)$, then
        \[
            \Prob(X=k)
            =
            \binom{n}{k}p^k(1-p)^{n-k},
            \qquad
            k=0,1,\dots,n.
        \]
        Its mean and variance are
        \[
            \Exp[X]=np,
            \qquad
            \Var(X)=np(1-p).
        \]
    \end{formula}

    \begin{formula}[Poisson Distribution]
        If $X\sim\Poi(\lambda)$, then
        \[
            \Prob(X=k)
            =
            e^{-\lambda}\frac{\lambda^k}{k!},
            \qquad
            k=0,1,2,\dots.
        \]
        Its mean and variance are
        \[
            \Exp[X]=\lambda,
            \qquad
            \Var(X)=\lambda.
        \]
    \end{formula}

    \begin{theorem}[Poisson Limit Theorem; Law of Small Numbers]
        If $n\to\infty$, $p\to 0$, and $np\to\lambda$, then
        \[
            \binom{n}{k}p^k(1-p)^{n-k}
            \to
            e^{-\lambda}\frac{\lambda^k}{k!}
        \]
        for each fixed $k\geq 0$.
    \end{theorem}

    \begin{remark}
        This theorem explains why the Poisson distribution models rare events:
        many trials, small success probability, and finite expected count.
    \end{remark}

    \begin{formula}[Normal Distribution; Gaussian Distribution]
        A random variable $X\sim\Normal(\mu,\sigma^2)$ has density
        \[
            f(x)
            =
            \frac{1}{\sigma\sqrt{2\pi}}
            \exp\left(
                -\frac{(x-\mu)^2}{2\sigma^2}
            \right).
        \]
        The standard normal distribution is $\Normal(0,1)$.
    \end{formula}

    \begin{formula}[Empirical Rule; 68--95--99.7 Rule]
        For a normal distribution, approximately
        \[
            68\%
        \]
        of the mass lies within one standard deviation of the mean,
        \[
            95\%
        \]
        lies within two standard deviations, and
        \[
            99.7\%
        \]
        lies within three standard deviations.
    \end{formula}

    \begin{theorem}[Weak Law of Large Numbers]
        Let $X_1,X_2,\dots$ be independent identically distributed random
        variables with mean $\mu$. Then, for every $\eps>0$,
        \[
            \Prob\left(
                \left|
                    \frac{X_1+\cdots+X_n}{n}-\mu
                \right|
                >
                \eps
            \right)
            \to 0
        \]
        as $n\to\infty$.
    \end{theorem}

    \begin{idea}
        The Law of Large Numbers says that the sample average stabilizes near
        the true mean when the number of trials becomes large.
    \end{idea}

    \begin{theorem}[Central Limit Theorem]
        Let $X_1,X_2,\dots$ be independent identically distributed random
        variables with mean $\mu$ and variance $\sigma^2>0$. Then
        \[
            \frac{X_1+\cdots+X_n-n\mu}{\sigma\sqrt n}
        \]
        converges in distribution to the standard normal distribution
        $\Normal(0,1)$.
    \end{theorem}

    \begin{remark}
        The Central Limit Theorem explains why the normal distribution appears
        so often in nature and statistics.
    \end{remark}

    \begin{theorem}[Markov's Inequality]
        If $X$ is a nonnegative random variable and $a>0$, then
        \[
            \Prob(X\geq a)
            \leq
            \frac{\Exp[X]}{a}.
        \]
    \end{theorem}

    \begin{theorem}[Chebyshev's Inequality]
        If $X$ has mean $\mu$ and variance $\sigma^2$, then for every $k>0$,
        \[
            \Prob(|X-\mu|\geq k\sigma)
            \leq
            \frac{1}{k^2}.
        \]
        Equivalently, for every $\eps>0$,
        \[
            \Prob(|X-\mu|\geq \eps)
            \leq
            \frac{\Var(X)}{\eps^2}.
        \]
    \end{theorem}

    \begin{remark}
        Markov's inequality uses only the mean, while Chebyshev's inequality
        uses the variance. Both are probability bounds that work without
        assuming a normal distribution.
    \end{remark}

    \begin{theorem}[Jensen's Inequality]
        If $\varphi$ is convex, then
        \[
            \varphi(\Exp[X])
            \leq
            \Exp[\varphi(X)].
        \]
        If $\varphi$ is concave, the inequality is reversed.
    \end{theorem}

    \begin{idea}
        Jensen's inequality says that, for a convex function, applying the
        function after averaging gives a value no larger than averaging after
        applying the function.
    \end{idea}

    \begin{formula}[Birthday Problem Approximation]
        In a group of $n$ people, assuming $365$ equally likely birthdays, the
        probability that at least two people share a birthday is approximately
        \[
            1-\exp\left(-\frac{n(n-1)}{730}\right).
        \]
        For $n=23$, this probability is already greater than $1/2$.
    \end{formula}

    \begin{remark}
        The birthday problem is famous because the answer is surprisingly
        small: only $23$ people are needed for a better-than-even chance of a
        shared birthday.
    \end{remark}

    \begin{formula}[Gambler's Ruin]
        Suppose a gambler starts with $i$ dollars and stops when reaching
        either $0$ dollars or $N$ dollars. If the probability of winning each
        round is $p$ and $q=1-p$, then the probability of reaching $N$ before
        $0$ is
        \[
            \frac{1-(q/p)^i}{1-(q/p)^N}
            \qquad
            (p\neq q).
        \]
        In the fair case $p=q=1/2$, the probability is
        \[
            \frac{i}{N}.
        \]
    \end{formula}

    \begin{remark}
        Gambler's ruin is a classic random walk result. It is a good example
        where the fair case has a very simple answer.
    \end{remark}

    \chapter{Geometry}

    This chapter collects quiz-friendly results from Euclidean synthetic
    geometry and analytic geometry. The emphasis is on named theorems, elegant
    geometric configurations, and formulas with clear memorable answers.

    \section*{Synthetic Geometry}

    \begin{definition}[Classical Triangle Centers]
        For a triangle $\triangle ABC$, the most important classical centers
        are:
        \[
            O=\text{circumcenter},
            \qquad
            I=\text{incenter},
            \qquad
            H=\text{orthocenter},
            \qquad
            G=\text{centroid}.
        \]
        Here $O$ is the center of the circumcircle, $I$ is the center of the
        incircle, $H$ is the intersection of the altitudes, and $G$ is the
        intersection of the medians.
    \end{definition}

    \begin{remark}
        Triangle centers are extremely common in olympiad-style geometry.
        The symbols $O,I,H,G$ are standard and should be recognized quickly.
    \end{remark}

    \begin{formula}[Extended Law of Sines]
        For a triangle $\triangle ABC$ with side lengths
        \[
            a=BC,\qquad b=CA,\qquad c=AB,
        \]
        and circumradius $R$, we have
        \[
            \frac{a}{\sin A}
            =
            \frac{b}{\sin B}
            =
            \frac{c}{\sin C}
            =
            2R.
        \]
    \end{formula}

    \begin{formula}[Heron's Formula]
        Let $\triangle ABC$ have side lengths $a,b,c$ and semiperimeter
        \[
            s=\frac{a+b+c}{2}.
        \]
        Then its area $K$ is
        \[
            K
            =
            \sqrt{s(s-a)(s-b)(s-c)}.
        \]
        Also,
        \[
            K=rs
            \qquad\text{and}\qquad
            K=\frac{abc}{4R},
        \]
        where $r$ is the inradius and $R$ is the circumradius.
    \end{formula}

    \begin{strategy}
        The formulas
        \[
            K=rs,
            \qquad
            K=\frac{abc}{4R}
        \]
        are often more useful than Heron's formula when incircles or
        circumcircles appear.
    \end{strategy}

    \begin{theorem}[Angle Bisector Theorem]
        In a triangle $\triangle ABC$, suppose the internal angle bisector of
        $\angle A$ meets $BC$ at $D$. Then
        \[
            \frac{BD}{DC}
            =
            \frac{AB}{AC}.
        \]
    \end{theorem}

    \begin{theorem}[Ceva's Theorem]
        Let $D,E,F$ lie on the lines $BC,CA,AB$, respectively. Then the cevians
        $AD,BE,CF$ are concurrent if and only if
        \[
            \frac{BD}{DC}
            \cdot
            \frac{CE}{EA}
            \cdot
            \frac{AF}{FB}
            =
            1,
        \]
        using directed lengths.
    \end{theorem}

    \begin{remark}
        Ceva's Theorem is the standard theorem for proving concurrence of three
        lines in a triangle.
    \end{remark}

    \begin{theorem}[Menelaus' Theorem]
        Let $D,E,F$ lie on the lines $BC,CA,AB$, respectively. Then the points
        $D,E,F$ are collinear if and only if
        \[
            \frac{BD}{DC}
            \cdot
            \frac{CE}{EA}
            \cdot
            \frac{AF}{FB}
            =
            -1,
        \]
        using directed lengths.
    \end{theorem}

    \begin{remark}
        Ceva is for concurrence, while Menelaus is for collinearity. This pair
        is one of the most important toolkits in olympiad geometry.
    \end{remark}

    \begin{theorem}[Euler Line Theorem]
        In any non-equilateral triangle, the circumcenter $O$, centroid $G$,
        and orthocenter $H$ are collinear. Moreover,
        \[
            OG:GH=1:2.
        \]
        The line passing through these points is called the Euler line.
    \end{theorem}

    \begin{theorem}[Nine-Point Circle Theorem; Euler Circle]
        In any triangle, the following nine points lie on one circle:
        the three side midpoints, the three feet of the altitudes, and the
        three midpoints between the orthocenter and the vertices. This circle
        is called the nine-point circle.
    \end{theorem}

    \begin{remark}
        The center of the nine-point circle is the midpoint of the segment
        joining the circumcenter $O$ and the orthocenter $H$.
    \end{remark}

    \begin{theorem}[Euler's Formula for the Incenter and Circumcenter]
        For a triangle with circumradius $R$, inradius $r$, circumcenter $O$,
        and incenter $I$, we have
        \[
            OI^2=R^2-2Rr.
        \]
        In particular,
        \[
            R\geq 2r.
        \]
    \end{theorem}

    \begin{remark}
        The inequality $R\geq 2r$ is called Euler's inequality for triangles.
        Equality holds exactly for an equilateral triangle.
    \end{remark}

    \begin{theorem}[Power of a Point Theorem]
        Let $P$ be a point and let a line through $P$ meet a circle at
        $A$ and $B$. The product
        \[
            PA\cdot PB
        \]
        is independent of the choice of the line through $P$. It is called the
        power of $P$ with respect to the circle.
    \end{theorem}

    \begin{formula}[Tangent--Secant Theorem]
        If $PT$ is tangent to a circle at $T$ and a secant through $P$ meets
        the circle at $A$ and $B$, then
        \[
            PT^2=PA\cdot PB.
        \]
    \end{formula}

    \begin{theorem}[Radical Axis Theorem]
        For two nonconcentric circles, the set of points having equal powers
        with respect to the two circles is a line. This line is called the
        radical axis of the two circles.
    \end{theorem}

    \begin{remark}
        Radical axes are powerful because they turn circle configurations into
        collinearity statements.
    \end{remark}

    \begin{theorem}[Ptolemy's Theorem]
        A quadrilateral $ABCD$ is cyclic if and only if
        \[
            AC\cdot BD
            =
            AB\cdot CD+BC\cdot DA
        \]
        for the appropriate cyclic order of the vertices.
    \end{theorem}

    \begin{remark}
        Ptolemy's Theorem is one of the most recognizable formulas for cyclic
        quadrilaterals.
    \end{remark}

    \begin{theorem}[Simson Line Theorem]
        Let $P$ be a point on the circumcircle of $\triangle ABC$. The feet of
        the perpendiculars from $P$ to the three lines $AB,BC,CA$ are collinear.
        This line is called the Simson line of $P$ with respect to
        $\triangle ABC$.
    \end{theorem}

    \begin{theorem}[Miquel's Theorem]
        In a complete quadrilateral, the circumcircles of the four triangles
        formed by choosing three of the four lines pass through a common point.
        This point is called the Miquel point.
    \end{theorem}

    \begin{remark}
        Miquel's Theorem is a standard source of hidden cyclic quadrilaterals
        in olympiad geometry.
    \end{remark}

    \begin{theorem}[Butterfly Theorem]
        Let $M$ be the midpoint of a chord $PQ$ of a circle. Through $M$, draw
        two other chords $AB$ and $CD$, and let $AD$ and $BC$ meet $PQ$ at
        $X$ and $Y$, respectively. Then $M$ is also the midpoint of $XY$.
    \end{theorem}

    \begin{theorem}[Morley's Trisector Theorem]
        In any triangle, the adjacent trisectors of the three angles intersect
        in three points that form an equilateral triangle.
    \end{theorem}

    \begin{remark}
        Morley's theorem is famous because the conclusion is surprisingly
        simple: angle trisectors always create an equilateral triangle.
    \end{remark}

    \begin{theorem}[Napoleon's Theorem]
        Construct equilateral triangles externally on the three sides of any
        triangle. Then the centers of these three equilateral triangles form an
        equilateral triangle.
    \end{theorem}

    \begin{theorem}[Classification of Platonic Solids]
        There are exactly five Platonic solids:
        \[
            \text{tetrahedron},\quad
            \text{cube},\quad
            \text{octahedron},\quad
            \text{dodecahedron},\quad
            \text{icosahedron}.
        \]
    \end{theorem}

    \begin{remark}
        The classification of Platonic solids is a classic result in Euclidean
        geometry and a very natural quiz fact.
    \end{remark}

    \section*{Analytic Geometry}

    \begin{formula}[Shoelace Formula; Gauss's Area Formula]
        For a polygon with vertices
        \[
            (x_1,y_1),(x_2,y_2),\dots,(x_n,y_n)
        \]
        in cyclic order, its area is
        \[
            \frac12
            \left|
            \sum_{i=1}^{n}
            (x_i y_{i+1}-x_{i+1}y_i)
            \right|,
        \]
        where
        \[
            (x_{n+1},y_{n+1})=(x_1,y_1).
        \]
    \end{formula}

    \begin{remark}
        The Shoelace Formula is one of the cleanest examples of analytic
        geometry: a geometric area becomes a determinant-like coordinate
        calculation.
    \end{remark}

    \begin{definition}[Conic Section]
        A conic section is a curve obtained by intersecting a plane with a
        double cone. The nondegenerate conics are:
        \[
            \text{ellipse},\qquad
            \text{parabola},\qquad
            \text{hyperbola}.
        \]
        Equivalently, a conic can be described using a focus, a directrix, and
        an eccentricity $e$.
    \end{definition}

    \begin{formula}[Focus--Directrix Definition of Conics]
        A conic is the set of points $P$ satisfying
        \[
            \frac{\operatorname{dist}(P,\text{focus})}
                 {\operatorname{dist}(P,\text{directrix})}
            =
            e.
        \]
        It is an ellipse if
        \[
            0<e<1,
        \]
        a parabola if
        \[
            e=1,
        \]
        and a hyperbola if
        \[
            e>1.
        \]
    \end{formula}

    \begin{formula}[Standard Forms of Conics]
        The standard form of an ellipse is
        \[
            \frac{x^2}{a^2}+\frac{y^2}{b^2}=1.
        \]
        The standard form of a hyperbola is
        \[
            \frac{x^2}{a^2}-\frac{y^2}{b^2}=1.
        \]
        The standard form of a parabola opening to the right is
        \[
            y^2=4px.
        \]
    \end{formula}

    \begin{formula}[Eccentricity of an Ellipse and a Hyperbola]
        For an ellipse
        \[
            \frac{x^2}{a^2}+\frac{y^2}{b^2}=1
            \qquad (a\geq b>0),
        \]
        the eccentricity is
        \[
            e=\frac{c}{a},
            \qquad
            c^2=a^2-b^2.
        \]
        For a hyperbola
        \[
            \frac{x^2}{a^2}-\frac{y^2}{b^2}=1,
        \]
        the eccentricity is
        \[
            e=\frac{c}{a},
            \qquad
            c^2=a^2+b^2.
        \]
    \end{formula}

    \begin{theorem}[Reflection Property of Conics]
        An ellipse reflects a ray from one focus to the other focus. A parabola
        reflects rays parallel to its axis through its focus. A hyperbola has
        the analogous reflection property involving its two foci and branches.
    \end{theorem}

    \begin{remark}
        The reflection property gives a geometric reason why conics appear in
        optics, satellite dishes, whispering galleries, and orbital mechanics.
    \end{remark}

    \begin{theorem}[Quadratic Classification of Conics]
        Consider a real quadratic equation
        \[
            Ax^2+Bxy+Cy^2+Dx+Ey+F=0.
        \]
        The discriminant
        \[
            B^2-4AC
        \]
        determines the type of the nondegenerate conic:
        \[
            B^2-4AC<0
            \quad\Rightarrow\quad
            \text{ellipse-type},
        \]
        \[
            B^2-4AC=0
            \quad\Rightarrow\quad
            \text{parabola-type},
        \]
        and
        \[
            B^2-4AC>0
            \quad\Rightarrow\quad
            \text{hyperbola-type}.
        \]
    \end{theorem}

    \begin{remark}
        This is the analytic-geometry counterpart of classifying conic sections
        by their shapes.
    \end{remark}

    \begin{formula}[Polar Equation of a Conic]
        With a focus at the pole, a conic has polar equation
        \[
            r
            =
            \frac{\ell}{1+e\cos\theta}
        \]
        after choosing a suitable axis. Here $e$ is the eccentricity and
        $\ell$ is the semi-latus rectum.
    \end{formula}

    \begin{remark}
        The same equation describes ellipses, parabolas, and hyperbolas
        depending on whether $e<1$, $e=1$, or $e>1$.
    \end{remark}

    \begin{definition}[Apollonius Circle]
        Given two fixed points $A$ and $B$ and a positive constant
        $k\neq 1$, the locus of points $P$ satisfying
        \[
            \frac{PA}{PB}=k
        \]
        is a circle. This circle is called an Apollonius circle.
    \end{definition}

    \begin{remark}
        Apollonius circles turn a ratio of distances into a concrete circle,
        which is a common hidden structure in geometry problems.
    \end{remark}

    \begin{theorem}[Descartes' Circle Theorem; Soddy Circle Theorem]
        Suppose four circles are mutually tangent, and let their curvatures be
        \[
            k_1,k_2,k_3,k_4,
        \]
        where curvature means the reciprocal of the radius, with sign convention
        for an enclosing circle. Then
        \[
            (k_1+k_2+k_3+k_4)^2
            =
            2(k_1^2+k_2^2+k_3^2+k_4^2).
        \]
    \end{theorem}

    \begin{remark}
        Descartes' Circle Theorem is a striking formula: tangency of four
        circles becomes one quadratic equation in their curvatures.
    \end{remark}

    \begin{theorem}[Pick's Theorem]
        Let $P$ be a simple polygon whose vertices lie on lattice points in the
        plane. If $I$ is the number of lattice points strictly inside $P$ and
        $B$ is the number of lattice points on the boundary of $P$, then
        \[
            \operatorname{Area}(P)
            =
            I+\frac{B}{2}-1.
        \]
    \end{theorem}

    \begin{remark}
        Pick's Theorem is a beautiful bridge between geometry and arithmetic:
        area is determined by counting lattice points.
    \end{remark}

    \chapter{Number Theory}

    This chapter collects famous number-theoretic results with strong quiz
    value: primes, congruences, reciprocity, sums of squares, and the
    distribution of primes.

    \begin{theorem}[Fundamental Theorem of Arithmetic]
        Every integer $n>1$ can be written as a product of prime numbers.
        Moreover, this factorization is unique up to the order of the factors.
        Equivalently,
        \[
            n=p_1^{\alpha_1}p_2^{\alpha_2}\cdots p_k^{\alpha_k},
        \]
        where $p_1,\dots,p_k$ are distinct primes and
        $\alpha_1,\dots,\alpha_k$ are positive integers, and this expression is
        unique up to reordering.
    \end{theorem}

    \begin{remark}
        This theorem is often called unique prime factorization. It is the
        reason prime numbers are treated as the basic building blocks of
        integers.
    \end{remark}

    \begin{theorem}[Euclid's Theorem; Infinitude of Primes]
        There are infinitely many prime numbers.
    \end{theorem}

    \begin{idea}
        If $p_1,\dots,p_n$ are finitely many primes, consider
        \[
            p_1p_2\cdots p_n+1.
        \]
        This number is not divisible by any prime on the list.
    \end{idea}

    \begin{theorem}[Bezout's Identity]
        Let $a,b\in\Z$, not both zero. Then there exist integers $x,y\in\Z$
        such that
        \[
            ax+by=\gcd(a,b).
        \]
        In particular, if $\gcd(a,b)=1$, then there exist integers $x,y\in\Z$
        such that
        \[
            ax+by=1.
        \]
    \end{theorem}

    \begin{remark}
        Bezout's identity is the algebraic heart of the Euclidean algorithm and
        modular inverses.
    \end{remark}

    \begin{theorem}[Chinese Remainder Theorem; CRT]
        Let $m_1,\dots,m_k$ be pairwise relatively prime positive integers.
        Then the system
        \[
            x\equiv a_1\pmod {m_1},
            \qquad
            x\equiv a_2\pmod {m_2},
            \qquad
            \dots,
            \qquad
            x\equiv a_k\pmod {m_k}
        \]
        has a solution. Moreover, the solution is unique modulo
        \[
            M=m_1m_2\cdots m_k.
        \]
    \end{theorem}

    \begin{strategy}
        CRT is the standard theorem for combining several congruence conditions
        into one congruence.
    \end{strategy}

    \begin{formula}[Euler's Phi Function Formula]
        If
        \[
            n=p_1^{\alpha_1}p_2^{\alpha_2}\cdots p_k^{\alpha_k}
        \]
        is the prime factorization of $n$, then
        \[
            \varphi(n)
            =
            n\prod_{i=1}^{k}\left(1-\frac{1}{p_i}\right).
        \]
        In particular, for a prime power,
        \[
            \varphi(p^\alpha)=p^\alpha-p^{\alpha-1}
            =
            p^{\alpha-1}(p-1).
        \]
    \end{formula}

    \begin{theorem}[Fermat's Little Theorem]
        Let $p$ be a prime number. If $p\nmid a$, then
        \[
            a^{p-1}\equiv 1\pmod p.
        \]
        Equivalently, for every integer $a$,
        \[
            a^p\equiv a\pmod p.
        \]
    \end{theorem}

    \begin{theorem}[Euler's Theorem]
        If $\gcd(a,n)=1$, then
        \[
            a^{\varphi(n)}\equiv 1\pmod n.
        \]
    \end{theorem}

    \begin{remark}
        Fermat's Little Theorem is the special case of Euler's Theorem when
        $n=p$ is prime, since $\varphi(p)=p-1$.
    \end{remark}

    \begin{theorem}[Wilson's Theorem]
        A positive integer $p>1$ is prime if and only if
        \[
            (p-1)!\equiv -1\pmod p.
        \]
    \end{theorem}

    \begin{remark}
        Wilson's theorem is famous because it gives a clean factorial
        characterization of primes, even though it is not efficient as a
        primality test.
    \end{remark}

    \begin{definition}[Legendre Symbol]
        Let $p$ be an odd prime. The Legendre symbol is defined by
        \[
            \left(\frac{a}{p}\right)
            =
            \begin{cases}
                1, & \text{if }a\text{ is a quadratic residue modulo }p, \\
                -1, & \text{if }a\text{ is a quadratic nonresidue modulo }p, \\
                0, & \text{if }p\mid a.
            \end{cases}
        \]
    \end{definition}

    \begin{theorem}[Euler's Criterion]
        Let $p$ be an odd prime. For any integer $a$,
        \[
            \left(\frac{a}{p}\right)
            \equiv
            a^{(p-1)/2}
            \pmod p.
        \]
        Equivalently, if $p\nmid a$, then
        \[
            a^{(p-1)/2}
            \equiv
            \begin{cases}
                1\pmod p, & \text{if }a\text{ is a quadratic residue modulo }p, \\
                -1\pmod p, & \text{if }a\text{ is a quadratic nonresidue modulo }p.
            \end{cases}
        \]
    \end{theorem}

    \begin{theorem}[Law of Quadratic Reciprocity; Gauss's Theorem of Quadratic Reciprocity]
        Let $p$ and $q$ be distinct odd primes. Then
        \[
            \left(\frac{p}{q}\right)
            \left(\frac{q}{p}\right)
            =
            (-1)^{\frac{p-1}{2}\frac{q-1}{2}}.
        \]
        Equivalently,
        \[
            \left(\frac{p}{q}\right)
            =
            \left(\frac{q}{p}\right)
        \]
        unless both $p$ and $q$ are congruent to $3$ modulo $4$.
    \end{theorem}

    \begin{remark}
        Quadratic reciprocity is one of the crown jewels of elementary number
        theory. Gauss called it the golden theorem.
    \end{remark}

    \begin{theorem}[Fermat's Two-Square Theorem]
        An odd prime $p$ can be written as a sum of two squares,
        \[
            p=a^2+b^2
        \]
        for some integers $a,b\in\Z$, if and only if
        \[
            p\equiv 1\pmod 4.
        \]
        Also,
        \[
            2=1^2+1^2.
        \]
    \end{theorem}

    \begin{theorem}[Lagrange's Four-Square Theorem]
        Every nonnegative integer can be written as a sum of four integer
        squares. That is, for every $n\in\Nzero$, there exist integers
        $a,b,c,d\in\Z$ such that
        \[
            n=a^2+b^2+c^2+d^2.
        \]
    \end{theorem}

    \begin{remark}
        The two-square theorem has a congruence condition. The four-square
        theorem has no exception at all.
    \end{remark}

    \begin{theorem}[Dirichlet's Theorem on Arithmetic Progressions]
        If $a$ and $d$ are relatively prime positive integers, then the
        arithmetic progression
        \[
            a,\ a+d,\ a+2d,\ a+3d,\dots
        \]
        contains infinitely many primes.
    \end{theorem}

    \begin{remark}
        Dirichlet's theorem is a major result connecting primes and arithmetic
        progressions. The condition $\gcd(a,d)=1$ is necessary.
    \end{remark}

    \begin{theorem}[Bertrand's Postulate; Chebyshev's Theorem]
        For every integer $n>1$, there exists a prime $p$ such that
        \[
            n<p<2n.
        \]
    \end{theorem}

    \begin{theorem}[Prime Number Theorem]
        Let $\pi(x)$ be the number of primes less than or equal to $x$. Then
        \[
            \pi(x)\sim \frac{x}{\log x}
            \qquad
            \text{as }x\to\infty.
        \]
    \end{theorem}

    \begin{idea}
        The Prime Number Theorem says that the ``probability'' that a large
        integer near $x$ is prime is roughly $1/\log x$. This is a mental model,
        not a literal probability statement.
    \end{idea}

    \begin{definition}[Riemann Zeta Function]
        For $\operatorname{Re}(s)>1$, the Riemann zeta function is defined by
        \[
            \zeta(s)
            =
            \sum_{n=1}^{\infty}\frac{1}{n^s}.
        \]
        It also has the Euler product
        \[
            \zeta(s)
            =
            \prod_{p\text{ prime}}
            \frac{1}{1-p^{-s}}.
        \]
    \end{definition}

    \begin{remark}
        The Euler product is the bridge between the zeta function and prime
        numbers. The Riemann Hypothesis asserts that the nontrivial zeros of
        $\zeta(s)$ all have real part $1/2$.
    \end{remark}

    \begin{theorem}[Euclid--Euler Theorem on Perfect Numbers]
        An even positive integer $N$ is perfect if and only if
        \[
            N=2^{p-1}(2^p-1),
        \]
        where $2^p-1$ is a Mersenne prime.
    \end{theorem}

    \begin{remark}
        A perfect number is a positive integer equal to the sum of its positive
        proper divisors. For example,
        \[
            6=1+2+3.
        \]
    \end{remark}

    \begin{theorem}[Fermat's Last Theorem]
        For every integer $n\geq 3$, there do not exist positive integers
        $a,b,c$ such that
        \[
            a^n+b^n=c^n.
        \]
    \end{theorem}

    \begin{remark}
        Fermat's Last Theorem was proved by Andrew Wiles. It is one of the most
        famous theorems in the history of mathematics.
    \end{remark}

    \chapter{Algebra}

    This chapter collects algebraic structures and named theorems that are likely
    to appear as terminology or quick-recognition questions. The emphasis is on
    groups, rings, fields, and polynomial equations, not on long abstract proofs.


    \begin{definition}[Group]
        A group is a set $G$ together with a binary operation satisfying closure,
        associativity, the existence of an identity element, and the existence of
        inverses. If the operation is commutative, then $G$ is called an abelian
        group.
    \end{definition}

    \begin{definition}[Cyclic Group]
        A group $G$ is cyclic if there exists an element $g\in G$ such that every
        element of $G$ is a power of $g$. In this case, we write
        \[
            G=\langle g\rangle.
        \]
    \end{definition}

    \begin{theorem}[Lagrange's Theorem]
        If $G$ is a finite group and $H\leq G$, then
        \[
            |H|\mid |G|.
        \]
        Equivalently, the order of every subgroup divides the order of the group.
    \end{theorem}

    \begin{corollary}[Order of an Element]
        If $G$ is a finite group and $g\in G$, then the order of $g$ divides
        $|G|$. In particular,
        \[
            g^{|G|}=e.
        \]
    \end{corollary}

    \begin{theorem}[Cauchy's Theorem]
        If $G$ is a finite group and a prime $p$ divides $|G|$, then $G$ contains
        an element of order $p$.
    \end{theorem}

    \begin{theorem}[First Isomorphism Theorem for Groups]
        If $\varphi:G\to H$ is a group homomorphism, then
        \[
            G/\Ker(\varphi)\cong \Image(\varphi).
        \]
    \end{theorem}

    \begin{theorem}[Cayley's Theorem]
        Every group is isomorphic to a subgroup of a symmetric group. In
        particular, every finite group of order $n$ is isomorphic to a subgroup
        of $S_n$.
    \end{theorem}

    \begin{theorem}[Sylow Theorems]
        Let $G$ be a finite group and suppose
        \[
            |G|=p^k m,
            \qquad
            p\nmid m,
        \]
        where $p$ is prime. Then $G$ has a subgroup of order $p^k$, called a
        Sylow $p$-subgroup. Moreover, all Sylow $p$-subgroups are conjugate, and
        the number $n_p$ of Sylow $p$-subgroups satisfies
        \[
            n_p\equiv 1\pmod p,
            \qquad
            n_p\mid m.
        \]
    \end{theorem}

    \begin{strategy}
        For Science Quiz, the most common finite-group quick check is divisibility:
        subgroup orders and element orders must divide the group order. When a
        prime divides $|G|$, Cauchy's theorem guarantees an element of that prime
        order.
    \end{strategy}


    \begin{definition}[Ring and Field]
        A ring is a set equipped with addition and multiplication satisfying the
        usual distributive laws. A field is a commutative ring in which every
        nonzero element has a multiplicative inverse.
    \end{definition}

    \begin{definition}[Integral Domain]
        An integral domain is a commutative ring with identity in which there are
        no zero divisors. In other words, if $ab=0$, then $a=0$ or $b=0$.
    \end{definition}

    \begin{theorem}[Polynomial Division Algorithm]
        Let $F$ be a field. If $f(x),g(x)\in F[x]$ and $g(x)\neq 0$, then there
        exist unique polynomials $q(x),r(x)\in F[x]$ such that
        \[
            f(x)=q(x)g(x)+r(x),
            \qquad
            r(x)=0\quad\text{or}\quad \deg r<\deg g.
        \]
    \end{theorem}

    \begin{theorem}[Factor Theorem]
        For a polynomial $f(x)\in F[x]$ and an element $a\in F$,
        \[
            f(a)=0
            \quad\Longleftrightarrow\quad
            x-a \text{ divides } f(x).
        \]
    \end{theorem}

    \begin{theorem}[Fundamental Theorem of Algebra]
        Every nonconstant complex polynomial has at least one complex root.
        Equivalently, every polynomial of degree $n\geq 1$ over $\C$ splits into
        $n$ linear factors over $\C$, counted with multiplicity.
    \end{theorem}

    \begin{definition}[Algebraically Closed Field]
        A field $F$ is algebraically closed if every nonconstant polynomial in
        $F[x]$ has a root in $F$.
    \end{definition}

    \begin{theorem}[Galois Correspondence]
        For a finite Galois extension $L/K$, there is an inclusion-reversing
        correspondence between intermediate fields $K\subseteq E\subseteq L$ and
        subgroups $H\leq \Gal(L/K)$, given by
        \[
            E\longmapsto \Gal(L/E),
            \qquad
            H\longmapsto L^H.
        \]
    \end{theorem}

    \begin{theorem}[Abel--Ruffini Theorem]
        There is no formula for the roots of a general polynomial of degree $5$
        or higher using only the field operations and radicals.
    \end{theorem}

    \begin{remark}
        The Abel--Ruffini theorem does not say that every quintic equation is
        unsolvable. It says that no universal radical formula exists for the
        general quintic.
    \end{remark}

    \chapter{Analysis}

    This chapter collects foundational results from real analysis that often
    appear behind rigorous calculus, limits, compactness, and convergence. The
    statements are kept theorem-oriented so that they can be recognized quickly
    in a quiz setting.


    \begin{definition}[Limit Superior and Limit Inferior]
        For a real sequence $(a_n)$, define
        \[
            \limsup_{n\to\infty} a_n
            =
            \lim_{n\to\infty}\sup_{k\geq n} a_k,
            \qquad
            \liminf_{n\to\infty} a_n
            =
            \lim_{n\to\infty}\inf_{k\geq n} a_k,
        \]
        whenever these extended real limits are allowed.
    \end{definition}

    \begin{theorem}[Monotone Convergence Theorem for Sequences]
        Every bounded monotone real sequence converges. More precisely, every
        increasing sequence bounded above converges to its supremum, and every
        decreasing sequence bounded below converges to its infimum.
    \end{theorem}

    \begin{theorem}[Bolzano--Weierstrass Theorem]
        Every bounded sequence in $\R^n$ has a convergent subsequence.
    \end{theorem}

    \begin{theorem}[Heine--Borel Theorem]
        A subset of $\R^n$ is compact if and only if it is closed and bounded.
    \end{theorem}

    \begin{theorem}[Uniform Continuity on Compact Sets]
        If $K\subseteq\R^n$ is compact and $f:K\to\R^m$ is continuous, then $f$
        is uniformly continuous on $K$.
    \end{theorem}

    \begin{theorem}[Weierstrass Extreme Value Theorem]
        If $K\subseteq\R^n$ is compact and $f:K\to\R$ is continuous, then $f$
        attains both a maximum and a minimum on $K$.
    \end{theorem}

    \begin{strategy}
        In analysis questions, compactness usually means that one may extract a
        convergent subsequence or that a continuous function attains extrema.
        In $\R^n$, the fast recognition phrase is ``closed and bounded.''
    \end{strategy}


    \begin{theorem}[Cauchy Criterion]
        A sequence $(a_n)$ in $\R$ or $\C$ converges if and only if it is Cauchy:
        for every $\eps>0$, there exists $N$ such that
        \[
            m,n\geq N
            \quad\Longrightarrow\quad
            |a_m-a_n|<\eps.
        \]
    \end{theorem}

    \begin{theorem}[Weierstrass M-Test]
        Let $(f_n)$ be a sequence of functions on a set $E$. If
        \[
            |f_n(x)|\leq M_n
            \quad\text{for all }x\in E
        \]
        and
        \[
            \sum_{n=1}^{\infty}M_n
        \]
        converges, then the series
        \[
            \sum_{n=1}^{\infty} f_n(x)
        \]
        converges uniformly on $E$.
    \end{theorem}

    \begin{theorem}[Uniform Limit Theorem]
        If $(f_n)$ is a sequence of continuous functions on a set $E$ and
        $f_n\to f$ uniformly on $E$, then $f$ is continuous on $E$.
    \end{theorem}

    \begin{theorem}[Term-by-Term Differentiation]
        Suppose $f_n$ are continuously differentiable on an interval and $f_n$
        converges at one point. If $f_n'$ converges uniformly, then $f_n$
        converges uniformly to a differentiable function $f$, and
        \[
            f'=\lim_{n\to\infty} f_n'.
        \]
    \end{theorem}

    \begin{theorem}[Arzelà--Ascoli Theorem]
        A uniformly bounded and equicontinuous family of real-valued continuous
        functions on a compact metric space has a uniformly convergent
        subsequence.
    \end{theorem}


    \begin{theorem}[Banach Fixed-Point Theorem; Contraction Mapping Theorem]
        Let $(X,d)$ be a complete metric space, and let $T:X\to X$ be a
        contraction. Then $T$ has a unique fixed point $x^\ast\in X$, and the
        iteration $x_{n+1}=T(x_n)$ converges to $x^\ast$ for every initial point
        $x_0\in X$.
    \end{theorem}

    \begin{theorem}[Baire Category Theorem]
        A complete metric space cannot be written as a countable union of nowhere
        dense closed sets. Equivalently, a countable intersection of open dense
        subsets of a complete metric space is dense.
    \end{theorem}

    \begin{theorem}[Lebesgue Dominated Convergence Theorem]
        If $f_n\to f$ almost everywhere and there exists an integrable function
        $g$ such that $|f_n|\leq g$ for all $n$, then
        \[
            \lim_{n\to\infty}\int f_n\dd\mu
            =
            \int f\dd\mu.
        \]
    \end{theorem}

    \begin{theorem}[Fubini's Theorem]
        Under suitable integrability assumptions, an iterated integral may be
        computed in either order:
        \[
            \int_X\left(\int_Y f(x,y)\dd y\right)\dd x
            =
            \int_Y\left(\int_X f(x,y)\dd x\right)\dd y.
        \]
    \end{theorem}

    \chapter{Topology}

    This chapter records topology terms and theorems that are famous enough to
    appear in mathematical culture and Science Quiz questions. The goal is not to
    develop topology from scratch, but to make the names and meanings recognizable.


    \begin{definition}[Topology]
        A topology on a set $X$ is a collection $\mathcal{T}$ of subsets of $X$
        whose elements are called open sets, satisfying:
        \begin{enumerate}
            \item[(i)] $\varnothing,X\in\mathcal{T}$;
            \item[(ii)] arbitrary unions of sets in $\mathcal{T}$ belong to
            $\mathcal{T}$;
            \item[(iii)] finite intersections of sets in $\mathcal{T}$ belong to
            $\mathcal{T}$.
        \end{enumerate}
    \end{definition}

    \begin{definition}[Compactness]
        A topological space $X$ is compact if every open cover of $X$ has a finite
        subcover.
    \end{definition}

    \begin{definition}[Connectedness]
        A topological space $X$ is connected if it cannot be written as the union
        of two disjoint nonempty open sets. Equivalently, the only subsets of
        $X$ that are both open and closed are $\varnothing$ and $X$.
    \end{definition}

    \begin{definition}[Homeomorphism]
        A homeomorphism is a bijective continuous map whose inverse is also
        continuous. Two spaces are homeomorphic if they are the same from the
        viewpoint of topology.
    \end{definition}

    \begin{theorem}[Brouwer Fixed-Point Theorem]
        Every continuous map from a closed ball in $\R^n$ to itself has at least
        one fixed point.
    \end{theorem}

    \begin{theorem}[Invariance of Domain]
        If $U\subseteq\R^n$ is open and $f:U\to\R^n$ is continuous and injective,
        then $f(U)$ is open and $f$ is a homeomorphism from $U$ onto $f(U)$.
    \end{theorem}


    \begin{definition}[Manifold]
        An $n$-dimensional topological manifold is a space that locally looks like
        $\R^n$. More precisely, each point has a neighborhood homeomorphic to an
        open subset of $\R^n$.
    \end{definition}

    \begin{definition}[Fundamental Group]
        The fundamental group $\pi_1(X,x_0)$ records loops in $X$ based at
        $x_0$, up to continuous deformation. It is one of the basic algebraic
        invariants of a topological space.
    \end{definition}

    \begin{fact}[Fundamental Groups of Familiar Spaces]
        The following basic examples are worth recognizing:
        \[
            \pi_1(S^1)\cong\Z,
            \qquad
            \pi_1(S^n)\cong 0\quad(n\geq 2),
            \qquad
            \pi_1(\mathbb{T}^2)\cong\Z^2.
        \]
    \end{fact}

    \begin{formula}[Euler Characteristic of a Surface]
        For a triangulated closed surface, the Euler characteristic is
        \[
            \chi=V-E+F.
        \]
        For an orientable closed surface of genus $g$,
        \[
            \chi=2-2g.
        \]
    \end{formula}

    \begin{theorem}[Classification of Closed Surfaces]
        Every connected closed surface is homeomorphic to either an orientable
        surface of genus $g\geq 0$ or a nonorientable surface obtained as a
        connected sum of $k\geq 1$ projective planes.
    \end{theorem}

    \begin{theorem}[Poincaré Conjecture]
        Every closed simply connected $3$-manifold is homeomorphic to the
        $3$-sphere $S^3$.
    \end{theorem}

    \begin{remark}
        The Poincaré conjecture was proved by Grigori Perelman using Ricci flow.
        In a quiz setting, the key association is usually
        \[
            \text{Poincaré conjecture} \longleftrightarrow
            \text{Perelman} \longleftrightarrow
            \text{Ricci flow}.
        \]
    \end{remark}

    \begin{theorem}[Thurston Geometrization Theorem]
        Every closed orientable $3$-manifold can be decomposed into pieces that
        admit one of the eight Thurston geometries.
    \end{theorem}

    \begin{fact}[The Eight Thurston Geometries]
        The eight model geometries are
        \[
            S^3,
            \quad \mathbb{E}^3,
            \quad \mathbb{H}^3,
            \quad S^2\times\R,
            \quad \mathbb{H}^2\times\R,
            \quad \widetilde{\mathrm{SL}_2(\R)},
            \quad \mathrm{Nil},
            \quad \mathrm{Sol}.
        \]
    \end{fact}

    \chapter{Miscellaneous Topics}


    \begin{formula}[Mercator Series for $\ln 2$]
        The alternating harmonic series satisfies
        \[
            \sum_{n=1}^{\infty}(-1)^{n+1}\frac{1}{n}
            =
            1-\frac12+\frac13-\frac14+\cdots
            =
            \ln 2.
        \]
    \end{formula}

    \begin{remark}
        This follows from the Mercator series
        \[
            \ln(1+x)
            =
            x-\frac{x^2}{2}+\frac{x^3}{3}-\frac{x^4}{4}+\cdots
        \]
        by substituting $x=1$.
    \end{remark}

    \begin{formula}[Basel Sum; Euler's Basel Formula]
        The Basel sum is
        \[
            \sum_{n=1}^{\infty}\frac{1}{n^2}
            =
            1+\frac{1}{2^2}+\frac{1}{3^2}+\cdots
            =
            \frac{\pi^2}{6}.
        \]
        Equivalently,
        \[
            \zeta(2)=\frac{\pi^2}{6}.
        \]
    \end{formula}

    \begin{remark}
        The problem of evaluating
        \[
            \sum_{n=1}^{\infty}\frac{1}{n^2}
        \]
        is called the Basel problem. Its famous solution was found by Euler.
    \end{remark}

    \begin{formula}[Odd Part of the Basel Sum]
        The sum of reciprocal squares over odd positive integers is
        \[
            \sum_{n=1}^{\infty}\frac{1}{(2n-1)^2}
            =
            1+\frac{1}{3^2}+\frac{1}{5^2}+\cdots
            =
            \frac{\pi^2}{8}.
        \]
    \end{formula}

    \begin{idea}
        This is obtained by subtracting the even part from the Basel sum:
        \[
            \sum_{n=1}^{\infty}\frac{1}{(2n-1)^2}
            =
            \sum_{n=1}^{\infty}\frac{1}{n^2}
            -
            \sum_{n=1}^{\infty}\frac{1}{(2n)^2}
            =
            \left(1-\frac14\right)\frac{\pi^2}{6}
            =
            \frac{\pi^2}{8}.
        \]
    \end{idea}

    \begin{formula}[Alternating Basel Series; Dirichlet Eta Value at $2$]
        The alternating Basel series satisfies
        \[
            \sum_{n=1}^{\infty}(-1)^{n+1}\frac{1}{n^2}
            =
            1-\frac{1}{2^2}+\frac{1}{3^2}-\frac{1}{4^2}+\cdots
            =
            \frac{\pi^2}{12}.
        \]
        Equivalently,
        \[
            \eta(2)=\frac{\pi^2}{12},
        \]
        where $\eta(s)$ is the Dirichlet eta function.
    \end{formula}

    \begin{idea}
        The positive terms are the odd reciprocal squares and the negative
        terms are the even reciprocal squares. Hence
        \[
            \sum_{n=1}^{\infty}(-1)^{n+1}\frac{1}{n^2}
            =
            \sum_{n=1}^{\infty}\frac{1}{(2n-1)^2}
            -
            \sum_{n=1}^{\infty}\frac{1}{(2n)^2}
            =
            \frac{\pi^2}{8}-\frac{\pi^2}{24}
            =
            \frac{\pi^2}{12}.
        \]
    \end{idea}

    \begin{formula}[Gregory--Leibniz Series; Leibniz Formula for $\pi$]
        The Gregory--Leibniz series is
        \[
            \sum_{n=1}^{\infty}(-1)^{n+1}\frac{1}{2n-1}
            =
            1-\frac13+\frac15-\frac17+\cdots
            =
            \frac{\pi}{4}.
        \]
        Equivalently,
        \[
            \pi
            =
            4\left(
                1-\frac13+\frac15-\frac17+\cdots
            \right).
        \]
    \end{formula}

    \begin{remark}
        This follows from the Taylor series
        \[
            \arctan x
            =
            x-\frac{x^3}{3}+\frac{x^5}{5}-\frac{x^7}{7}+\cdots
        \]
        by substituting $x=1$, since $\arctan 1=\pi/4$.
    \end{remark}

    \begin{strategy}
        These values are useful because Science Quiz problems often hide them
        behind a simple transformation. For example,
        \[
            \frac12-\frac14+\frac16-\frac18+\cdots
            =
            \frac12
            \left(
                1-\frac12+\frac13-\frac14+\cdots
            \right)
            =
            \frac{\ln 2}{2}.
        \]
    \end{strategy}

    \part{Facts to Memorize}

    \chapter{ICM and the Fields Medal}

    The International Congress of Mathematicians, commonly abbreviated as ICM,
    is the most important international congress in mathematics. The Fields
    Medal is awarded by the International Mathematical Union at the ICM every
    four years. A Fields Medalist must satisfy the age condition: the candidate's
    $40$th birthday must not occur before January $1$ of the year of the
    Congress.

    \vspace{1em}

    {\small
    \renewcommand{\arraystretch}{1.18}
    \begin{longtable}{@{}L{0.07\textwidth}L{0.15\textwidth}L{0.51\textwidth}L{0.18\textwidth}@{}}
        \toprule
        Year & ICM location & Medalists and areas & Notes \\
        \midrule
        \endfirsthead

        \toprule
        Year & ICM location & Medalists and areas & Notes \\
        \midrule
        \endhead

        \midrule
        \multicolumn{4}{r}{Continued on the next page.} \\
        \endfoot

        \bottomrule
        \endlastfoot

        1936
        & Oslo, Norway
        & Lars Ahlfors --- complex analysis, Riemann surfaces;
          Jesse Douglas --- minimal surfaces, Plateau problem
        & First Fields Medals. \\

        1950
        & Cambridge, USA
        & Laurent Schwartz --- distribution theory;
          Atle Selberg --- analytic number theory, sieve methods, zeta functions
        & First awards after World War II. \\

        1954
        & Amsterdam, Netherlands
        & Kunihiko Kodaira --- complex geometry, algebraic geometry;
          Jean-Pierre Serre --- algebraic topology, sheaf theory
        & Serre is the youngest Fields Medalist. \\

        1958
        & Edinburgh, UK
        & Klaus Roth --- number theory, Diophantine approximation;
          René Thom --- cobordism, differential topology
        &  \\

        1962
        & Stockholm, Sweden
        & Lars Hörmander --- partial differential equations;
          John Milnor --- differential topology, exotic spheres
        &  \\

        1966
        & Moscow, USSR
        & Michael Atiyah --- K-theory, index theorem;
          Paul Cohen --- set theory, forcing, continuum hypothesis;
          Alexander Grothendieck --- algebraic geometry, homological algebra;
          Stephen Smale --- differential topology, generalized Poincaré conjecture
        & First year with four medalists. \\

        1970
        & Nice, France
        & Alan Baker --- transcendental number theory;
          Heisuke Hironaka --- algebraic geometry, resolution of singularities;
          Serge Novikov --- topology, differential geometry;
          John G. Thompson --- finite group theory
        &  \\

        1974
        & Vancouver, Canada
        & Enrico Bombieri --- number theory, complex analysis, minimal surfaces;
          David Mumford --- algebraic geometry, moduli spaces
        &  \\

        1978
        & Helsinki, Finland
        & Pierre Deligne --- algebraic geometry, Weil conjectures;
          Charles Fefferman --- analysis, several complex variables;
          Grigory Margulis --- Lie groups, ergodic theory, discrete groups;
          Daniel Quillen --- algebraic K-theory
        &  \\

        1982
        & Warsaw, Poland
        & Alain Connes --- operator algebras, noncommutative geometry;
          William Thurston --- low-dimensional topology, $3$-manifolds;
          Shing-Tung Yau --- differential geometry, PDE, Calabi conjecture
        & ICM Warsaw was held in 1983. \\

        1986
        & Berkeley, USA
        & Simon Donaldson --- four-manifold topology, gauge theory;
          Gerd Faltings --- arithmetic algebraic geometry, Mordell conjecture;
          Michael Freedman --- four-manifold topology, four-dimensional Poincaré conjecture
        &  \\

        1990
        & Kyoto, Japan
        & Vladimir Drinfeld --- quantum groups, Langlands program;
          Vaughan Jones --- operator algebras, knot theory;
          Shigefumi Mori --- algebraic geometry, minimal model program;
          Edward Witten --- mathematical physics, topology, quantum field theory
        & Witten was the first physicist to receive the Fields Medal. \\

        1994
        & Zurich, Switzerland
        & Jean Bourgain --- analysis, Banach spaces, harmonic analysis;
          Pierre-Louis Lions --- partial differential equations;
          Jean-Christophe Yoccoz --- dynamical systems;
          Efim Zelmanov --- algebra, restricted Burnside problem
        &  \\

        1998
        & Berlin, Germany
        & Richard Borcherds --- algebra, automorphic forms, moonshine;
          W. Timothy Gowers --- functional analysis, combinatorics;
          Maxim Kontsevich --- algebraic geometry, topology, mathematical physics;
          Curtis T. McMullen --- complex dynamics, hyperbolic geometry
        & Andrew Wiles received a special IMU silver plaque. \\

        2002
        & Beijing, China
        & Laurent Lafforgue --- Langlands correspondence over function fields;
          Vladimir Voevodsky --- motivic cohomology, $\mathbb{A}^1$-homotopy theory
        &  \\

        2006
        & Madrid, Spain
        & Andrei Okounkov --- probability, representation theory, algebraic geometry;
          Grigori Perelman --- Ricci flow, Poincaré conjecture;
          Terence Tao --- PDE, harmonic analysis, additive combinatorics;
          Wendelin Werner --- probability, SLE, Brownian motion
        & Perelman declined the medal. \\

        2010
        & Hyderabad, India
        & Elon Lindenstrauss --- ergodic theory, number theory;
          Ngo Bao Chau --- automorphic forms, Fundamental Lemma;
          Stanislav Smirnov --- probability, statistical physics, conformal invariance;
          Cédric Villani --- kinetic theory, Boltzmann equation, optimal transport
        &  \\

        2014
        & Seoul, South Korea
        & Artur Avila --- dynamical systems;
          Manjul Bhargava --- number theory, geometry of numbers;
          Martin Hairer --- stochastic partial differential equations;
          Maryam Mirzakhani --- Riemann surfaces, moduli spaces, dynamics
        & Mirzakhani was the first woman Fields Medalist. \\

        2018
        & Rio de Janeiro, Brazil
        & Caucher Birkar --- algebraic geometry, Fano varieties, minimal model program;
          Alessio Figalli --- optimal transport, PDE, metric geometry;
          Peter Scholze --- arithmetic algebraic geometry, $p$-adic geometry;
          Akshay Venkatesh --- number theory, homogeneous dynamics, representation theory
        &  \\

        2022
        & Helsinki, Finland
        & Hugo Duminil-Copin --- probability, statistical physics, phase transitions;
          June Huh --- Hodge theory, combinatorics, matroids;
          James Maynard --- analytic number theory, prime numbers;
          Maryna Viazovska --- sphere packing, $E_8$ lattice, Fourier analysis
        & ICM 2022 was virtual; the award ceremony was held in Helsinki.
          Viazovska was the second woman Fields Medalist. \\
    \end{longtable}
    }

    \section*{Essential ICM and Fields Medal Landmarks}

    \begin{itemize}
        \item ICM stands for International Congress of Mathematicians.
        \item IMU stands for International Mathematical Union.
        \item The Fields Medal is awarded every four years at the ICM.
        \item The Fields Medal is named after John Charles Fields.
        \item The Fields Medal was first awarded in 1936.
        \item The first two Fields Medalists were Lars Ahlfors and Jesse Douglas.
        \item A Fields Medalist must satisfy the age condition: the candidate's
              $40$th birthday must not occur before January $1$ of the year of
              the Congress.
        \item Jean-Pierre Serre is the youngest Fields Medalist.
        \item Edward Witten was the first physicist to receive the Fields Medal.
        \item Andrew Wiles received a special IMU silver plaque in 1998 for his
              proof of Fermat's Last Theorem.
        \item Grigori Perelman declined the Fields Medal in 2006 after proving
              the Poincaré conjecture.
        \item Maryam Mirzakhani was the first woman Fields Medalist.
        \item Maryna Viazovska was the second woman Fields Medalist.
        \item The 2026 ICM is scheduled to be held in Philadelphia, USA.
    \end{itemize}

    \part{Past Problems}

    \begin{problem}{2018.1}
        The following people are this year's Fields Medalists. Excluding choice (5), who is the second oldest among them?
        \[
            \begin{array}{ll}
                (1) & \text{Caucher Birkar} \\
                (2) & \text{Alessio Figalli} \\
                (3) & \text{Peter Scholze} \\
                (4) & \text{Akshay Venkatesh} \\
                (5) & \text{The person who stole Birkar's medal}
            \end{array}
        \]
    \end{problem}

    \begin{solution}
        The 2018 Fields Medalists were
        \[
            \begin{array}{lll}
                (1) & \text{Caucher Birkar} & \text{born in 1978}, \\
                (2) & \text{Alessio Figalli} & \text{born in 1984}, \\
                (3) & \text{Peter Scholze} & \text{born in 1987}, \\
                (4) & \text{Akshay Venkatesh} & \text{born in 1981}.
            \end{array}
        \]
        Among these four, the second oldest is Akshay Venkatesh. Therefore, the answer is
        \[
            \boxed{(4)}.
        \]

        Choice (5) is a joke referring to the incident in which Caucher Birkar's Fields Medal was stolen shortly after the 2018 Fields Medal ceremony. Birkar later received a replacement medal.
    \end{solution}

    \begin{problem}{2018.2}
        According to William Thurston's geometrization conjecture, a $3$-dimensional manifold can be cut along spheres and tori so that each resulting piece has one of $n$ possible geometric structures. What is $n$?
    \end{problem}

    \begin{solution}
        The $3$-dimensional geometric structures appearing in Thurston's geometrization conjecture are the eight Thurston geometries:
        \[
            S^3,\quad
            \Exp^3,\quad
            \mathbb{H}^3,\quad
            S^2 \times \R,\quad
            \mathbb{H}^2 \times \R,\quad
            \widetilde{\mathrm{SL}_2(\R)},\quad
            \mathrm{Nil},\quad
            \mathrm{Sol}.
        \]
        Hence $n=8$, so the answer is
        \[
            \boxed{8}.
        \]
    \end{solution}

    \begin{problem}{2018.3}
        Find the sum of the following infinite series:
        \[
            \frac{1}{2}-\frac{1}{4}+\frac{1}{6}-\frac{1}{8}
            +\frac{1}{10}-\frac{1}{12}+\cdots.
        \]
    \end{problem}

    \begin{solution}
        The given series is one half of the alternating harmonic series:
        \[
            \frac{1}{2}-\frac{1}{4}+\frac{1}{6}-\frac{1}{8}+\cdots
            =
            \frac{1}{2}
            \left(
                1-\frac{1}{2}+\frac{1}{3}-\frac{1}{4}+\cdots
            \right).
        \]
        Recall the Taylor expansion
        \[
            \ln(1+x)
            =
            x-\frac{x^2}{2}+\frac{x^3}{3}-\frac{x^4}{4}+\cdots.
        \]
        Substituting $x=1$, we obtain
        \[
            \ln 2
            =
            1-\frac{1}{2}+\frac{1}{3}-\frac{1}{4}+\cdots.
        \]
        Therefore,
        \[
            \frac{1}{2}-\frac{1}{4}+\frac{1}{6}-\frac{1}{8}+\cdots
            =
            \frac{\ln 2}{2}.
        \]
        Thus the answer is $\boxed{\dfrac{\ln 2}{2}}$.
    \end{solution}

    \begin{problem}{2018.4}
        Among the following five Fields Medalists, which one does not have a Korean doctoral student?
        \[
            \begin{array}{ll}
                (1) & \text{William Thurston} \\
                (2) & \text{Maxim Kontsevich} \\
                (3) & \text{Grigory Margulis} \\
                (4) & \text{Terence Tao} \\
                (5) & \text{Efim Zelmanov}
            \end{array}
        \]
    \end{problem}

    \begin{solution}
        The intended answer is Maxim Kontsevich.

        This problem is not asking for a theorem but for advisor--student data
        about Fields Medalists and Korean mathematicians. In a Science Quiz
        setting, it is best treated as a memorization-style question: among the
        five listed Fields Medalists, the exceptional name in the intended answer
        key is
        \[
            \boxed{(2)\ \text{Maxim Kontsevich}}.
        \]
    \end{solution}

    \begin{problem}{2020.1}
        The International Congress of Mathematicians, abbreviated as ICM, is a congress for mathematicians from around the world. It is organized every four years by the International Mathematical Union and may be thought of as a kind of Olympic Games for mathematicians. The 2014 ICM was held in South Korea, and the 2018 ICM was held in Brazil. As of 2019, in which country was the 2022 ICM scheduled to be held?
    \end{problem}

    \begin{solution}
        ICM stands for the International Congress of Mathematicians. It is the main international congress in mathematics and is organized every four years by the International Mathematical Union.

        The 2014 ICM was held in Seoul, South Korea, and the 2018 ICM was held in Rio de Janeiro, Brazil. At the time this problem was asked in 2019, the 2022 ICM was scheduled to be held in Saint Petersburg, Russia. Therefore, the answer is $\boxed{\text{Russia}}$.

        In reality, after Russia's invasion of Ukraine in 2022, the 2022 ICM was not held in Saint Petersburg. It was eventually held as a virtual event, with the award ceremony taking place in Helsinki, Finland. Thus, the answer to this problem should be understood according to the planned host country at the time of the 2019 competition.
    \end{solution}

    \begin{problem}{2020.2}
        A fair six-sided die has the numbers $-2,0,2,3,4,5$ written on its faces. The die is rolled $10$ times, and $X$ is defined to be the product of the $10$ numbers obtained. What is the expected value of $X$?
    \end{problem}

    \begin{solution}
        Let $Y$ be the number obtained from one roll of the die. Since all six faces are equally likely,
        \[
            \Exp[Y]
            =
            \frac{-2+0+2+3+4+5}{6}
            =
            2.
        \]
        Now let $Y_1,Y_2,\dots,Y_{10}$ be the numbers obtained from the $10$ rolls. Then
        \[
            X=Y_1Y_2\cdots Y_{10}.
        \]
        Since the rolls are independent, the random variables $Y_1,Y_2,\dots,Y_{10}$ are independent. Therefore,
        \[
            \Exp[X]
            =
            \Exp[Y_1Y_2\cdots Y_{10}]
            =
            \Exp[Y_1]\Exp[Y_2]\cdots\Exp[Y_{10}].
        \]
        All the $Y_i$ have the same distribution as $Y$, so
        \[
            \Exp[X]
            =
            \left(\Exp[Y]\right)^{10}
            =
            2^{10}
            =
            1024.
        \]
        Thus the answer is $\boxed{1024}$.
    \end{solution}

    \begin{problem}{2020.3}
        Which of the following Fields Medalists was not serving as a professor at the time?
        \[
            \begin{array}{ll}
                (1) & \text{Cédric Villani} \\
                (2) & \text{Elon Lindenstrauss} \\
                (3) & \text{Peter Scholze} \\
                (4) & \text{Martin Hairer} \\
                (5) & \text{Terence Tao}
            \end{array}
        \]
    \end{problem}

    \begin{solution}
        At the time this problem was asked, Cédric Villani was active as a politician in France rather than serving as a professor.

        The other choices were mathematicians holding professorial positions. Therefore, the answer is $\boxed{(1)\ \text{Cédric Villani}}$.
    \end{solution}

    \begin{problem}{2020.4}
        In a three-dimensional gravitational field, when an object moves from a point $\mathbf{a}$ to another point $\mathbf{b}$, the change in its potential energy is independent of the path taken from $\mathbf{a}$ to $\mathbf{b}$. Which of the following mathematical facts is least directly related to this physical phenomenon?
        \[
            \begin{array}{ll}
                (1) & \text{A curl-free vector field is the gradient of some potential function.} \\
                (2) & \text{Stokes' theorem} \\
                (3) & \text{The Laplacian of } \dfrac{1}{|\mathbf{x}|} \text{ is a delta function, up to a constant factor.} \\
                (4) & \text{The curl of a gradient is identically zero.}
            \end{array}
        \]
    \end{problem}

    \begin{solution}
        The statement that the change in potential energy is independent of the path means that the gravitational field is a conservative vector field.

        Let $\mathbf{F}$ be the force field. The change in potential energy from $\mathbf{a}$ to $\mathbf{b}$ is given by
        \[
            \Delta U
            =
            -\int_{\mathbf{a}}^{\mathbf{b}} \mathbf{F}\cdot\dd\mathbf{r}.
        \]
        For this quantity to be independent of the path, the force field should be expressible as the negative gradient of a potential function:
        \[
            \mathbf{F}=-\nabla U.
        \]
        Thus choice (1) is directly related to the phenomenon.

        Also, the curl of a gradient is always zero:
        \[
            \nabla \times \nabla U=\mathbf{0}.
        \]
        Hence choice (4) is also directly related. Stokes' theorem says that
        \[
            \oint_{\partial S}\mathbf{F}\cdot\dd\mathbf{r}
            =
            \iint_S (\nabla\times \mathbf{F})\cdot \mathbf{n}\dd S.
        \]
        Therefore, if $\nabla\times\mathbf{F}=\mathbf{0}$, then the line integral around a closed curve is zero, which is another way to understand path independence. Thus choice (2) is also related.

        On the other hand,
        \[
            \Delta\left(\frac{1}{|\mathbf{x}|}\right)
            =
            -4\pi\delta(\mathbf{x})
        \]
        is related to the fundamental solution of the Laplacian and to Poisson's equation for gravitational potential. Although it is related to gravity, it is not the main mathematical reason why the potential energy change is independent of path. Therefore, the answer is $\boxed{(3)}$.
    \end{solution}

    \begin{problem}{2020.5}
        Among the following mathematics-related academic institutions, which one is actually a degree-granting school?
        \[
            \begin{array}{ll}
                (1) & \text{Institut des Hautes Études Scientifiques, IHES} \\
                (2) & \text{École normale supérieure} \\
                (3) & \text{Mathematical Sciences Research Institute, MSRI} \\
                (4) & \text{Institute for Advanced Study, IAS} \\
                (5) & \text{Collège de France}
            \end{array}
        \]
    \end{problem}

    \begin{solution}
        École normale supérieure is one of France's most prestigious higher education institutions. It operates educational programs and is naturally regarded as a degree-granting school.

        By contrast, IHES, MSRI, and IAS are research institutes rather than degree-granting schools. Collège de France is also a famous academic institution that offers public lectures, but it does not grant degrees.

        Therefore, the answer is $\boxed{(2)\ \text{École normale supérieure}}$.
    \end{solution}

    \begin{problem}{2021.1}
        Which is larger, $A$ or $B$, where
        \[
            A
            =
            1+\frac{1}{1+\frac{1}{1-\frac{1}{1+4}}},
            \qquad
            B
            =
            1+\frac{1}{1+\frac{1}{1+\frac{1}{1+3}}}?
        \]
    \end{problem}

    \begin{solution}
        First, we compute $A$:
        \[
            A
            =
            1+\frac{1}{1+\frac{1}{1-\frac{1}{5}}}
            =
            1+\frac{1}{1+\frac{1}{\frac{4}{5}}}
            =
            1+\frac{1}{1+\frac{5}{4}}
            =
            1+\frac{1}{\frac{9}{4}}
            =
            1+\frac{4}{9}
            =
            \frac{13}{9}.
        \]
        Next, we compute $B$:
        \[
            B
            =
            1+\frac{1}{1+\frac{1}{1+\frac{1}{4}}}
            =
            1+\frac{1}{1+\frac{1}{\frac{5}{4}}}
            =
            1+\frac{1}{1+\frac{4}{5}}
            =
            1+\frac{1}{\frac{9}{5}}
            =
            1+\frac{5}{9}
            =
            \frac{14}{9}.
        \]
        Since
        \[
            \frac{13}{9}<\frac{14}{9},
        \]
        we have $A<B$. Therefore, the answer is $\boxed{B}$.
    \end{solution}

    \begin{problem}{2021.2}
        Let
        \[
            A(n)
            =
            \frac{1}{\sqrt{5}}
            \left(
                \frac{1+\sqrt{5}}{2}
            \right)^{n+1}.
        \]
        What is the integer closest to $A(7)$?
    \end{problem}

    \begin{solution}
        Recall Binet's formula for the Fibonacci sequence:
        \[
            F_n
            =
            \frac{\varphi^n-\psi^n}{\sqrt{5}},
            \qquad
            \varphi=\frac{1+\sqrt{5}}{2},
            \qquad
            \psi=\frac{1-\sqrt{5}}{2}.
        \]
        Here
        \[
            A(7)
            =
            \frac{\varphi^8}{\sqrt{5}}.
        \]
        By Binet's formula,
        \[
            F_8
            =
            \frac{\varphi^8-\psi^8}{\sqrt{5}}.
        \]
        Since $|\psi|<1$, the term $\dfrac{\psi^8}{\sqrt{5}}$ is very small. Thus
        \[
            A(7)
            =
            \frac{\varphi^8}{\sqrt{5}}
            =
            F_8+\frac{\psi^8}{\sqrt{5}}
        \]
        is very close to $F_8$.

        The Fibonacci numbers are
        \[
            F_0=0,\quad F_1=1,\quad F_2=1,\quad F_3=2,\quad F_4=3,\quad
            F_5=5,\quad F_6=8,\quad F_7=13,\quad F_8=21.
        \]
        Therefore, the integer closest to $A(7)$ is $\boxed{21}$.
    \end{solution}

    \begin{problem}{2021.3}
        Choose all groups among the following that are not finitely generated:
        \[
            \begin{array}{ll}
                (1) & \text{The fundamental group of an open unit disk with finitely many punctures.} \\
                (2) & \mathrm{SL}_2(\Z), \text{ the group of } 2\times 2 \text{ integer matrices with determinant } 1. \\
                (3) & \C^{\times}, \text{ the group of nonzero complex numbers under multiplication.} \\
                (4) & \text{The Monster group, the largest sporadic finite simple group.}
            \end{array}
        \]
    \end{problem}

    \begin{solution}
        First, consider an open unit disk with finitely many punctures. If the removed points are $p_1,\dots,p_n$, then the fundamental group is isomorphic to the free group on $n$ generators:
        \[
            \pi_1\left(D\setminus\{p_1,\dots,p_n\}\right)\cong F_n.
        \]
        Hence the group in choice (1) is finitely generated.

        The group $\mathrm{SL}_2(\Z)$ is also finitely generated. For example, it is generated by the two matrices
        \[
            S=
            \begin{pmatrix}
                0 & -1 \\
                1 & 0
            \end{pmatrix},
            \qquad
            T=
            \begin{pmatrix}
                1 & 1 \\
                0 & 1
            \end{pmatrix}.
        \]
        Thus choice (2) is finitely generated.

        On the other hand, $\C^{\times}$ is an uncountable abelian group. Every finitely generated group is countable, so $\C^{\times}$ cannot be finitely generated. Therefore, choice (3) is not finitely generated.

        Finally, the Monster group is a finite group. Every finite group is finitely generated, since one may simply take all of its elements as generators. Hence choice (4) is finitely generated.

        Therefore, the only group in the list that is not finitely generated is $\boxed{(3)\ \C^{\times}}$.
    \end{solution}

    \begin{problem}{2021.4}
        Let
        \[
            \mathbf{A}
            =
            \begin{pmatrix}
                1 & 2 & 3 \\
                1 & 2 & 3 \\
                1 & 2 & 3
            \end{pmatrix}.
        \]
        Find all eigenvalues of $\mathbf{A}$ together with their multiplicities. Also determine whether $\mathbf{A}$ is diagonalizable.
    \end{problem}

    \begin{solution}
        We can write $\mathbf{A}$ as
        \[
            \mathbf{A}
            =
            \begin{pmatrix}
                1 \\
                1 \\
                1
            \end{pmatrix}
            \begin{pmatrix}
                1 & 2 & 3
            \end{pmatrix}.
        \]
        Thus $\mathbf{A}$ is a rank-one matrix.

        First, let
        \[
            \mathbf{v}
            =
            \begin{pmatrix}
                1 \\
                1 \\
                1
            \end{pmatrix}.
        \]
        Then
        \[
            \mathbf{A}\mathbf{v}
            =
            \begin{pmatrix}
                6 \\
                6 \\
                6
            \end{pmatrix}
            =
            6\mathbf{v}.
        \]
        Hence $6$ is an eigenvalue of $\mathbf{A}$.

        Since $\mathbf{A}$ has rank $1$, its nullity is $2$. Therefore, the eigenspace corresponding to the eigenvalue $0$ has dimension $2$. Indeed,
        \[
            \mathbf{A}
            \begin{pmatrix}
                x \\
                y \\
                z
            \end{pmatrix}
            =
            \begin{pmatrix}
                x+2y+3z \\
                x+2y+3z \\
                x+2y+3z
            \end{pmatrix},
        \]
        so
        \[
            E_0
            =
            \left\{
                \begin{pmatrix}
                    x \\
                    y \\
                    z
                \end{pmatrix}
                \in \R^3
                :
                x+2y+3z=0
            \right\}.
        \]
        Thus the geometric multiplicity of $0$ is $2$.

        Now
        \[
            \operatorname{tr}(\mathbf{A})=1+2+3=6.
        \]
        Since the sum of the eigenvalues, counted with algebraic multiplicity, is equal to the trace, the eigenvalues are
        \[
            6,\quad 0,\quad 0.
        \]
        Hence the algebraic multiplicities are
        \[
            6:1,
            \qquad
            0:2.
        \]

        Finally, the eigenspace for $6$ has dimension $1$, and the eigenspace for $0$ has dimension $2$. The total dimension of the eigenspaces is therefore
        \[
            1+2=3.
        \]
        Hence there exists a basis of $\R^3$ consisting of eigenvectors of $\mathbf{A}$. Therefore, $\mathbf{A}$ is diagonalizable.

        Thus the eigenvalues are $\boxed{6,0,0}$, with algebraic multiplicities $1$ and $2$ respectively, and $\boxed{\mathbf{A}\text{ is diagonalizable}}$.
    \end{solution}

    \begin{problem}{2022.1}
        Which of the following numbers is the largest?
        \[
            \begin{array}{ll}
                (1) & e^\pi \\
                (2) & \pi^e \\
                (3) & 7\pi \\
                (4) & \binom{7}{2} \\
                (5) & 4!
            \end{array}
        \]
    \end{problem}

    \begin{solution}
        We compare the values directly:
        \[
            e^\pi \approx 23.1407,
            \qquad
            \pi^e \approx 22.4592,
            \qquad
            7\pi \approx 21.9911.
        \]
        Also,
        \[
            \binom{7}{2}
            =
            \frac{7\cdot 6}{2}
            =
            21,
            \qquad
            4!
            =
            24.
        \]
        Therefore, the largest number among the choices is $\boxed{(5)\ 4!}$.
    \end{solution}

    \begin{problem}{2022.2}
        Three people store a jointly owned treasure in a safe and want to lock it. They want no single person to be able to open the safe alone, but any two or more people should always be able to open it together. This can be done by putting $n$ locks on the safe, making $m$ copies of the key for each lock, and distributing the keys appropriately among the three people. Find $n$ and $m$ so that the number of locks $n$ is minimized.
    \end{problem}

    \begin{solution}
        Let the three people be $A$, $B$, and $C$.

        Since any two people should be able to open the safe, the key to each lock must be given to at least two people. If a certain lock had its key given to only one person, then the other two people would not be able to open that lock together.

        Thus, for each lock, exactly one of the three people should not receive the key. On the other hand, no single person should be able to open the safe alone. Therefore, each person must be missing the key to at least one lock.

        Hence we need at least three locks: one lock whose key is not given to $A$, one lock whose key is not given to $B$, and one lock whose key is not given to $C$. This lower bound is achievable by distributing the keys as follows:
        \[
            \begin{array}{c|ccc}
                & A & B & C \\
                \hline
                L_A & \times & \checkmark & \checkmark \\
                L_B & \checkmark & \times & \checkmark \\
                L_C & \checkmark & \checkmark & \times
            \end{array}
        \]
        That is, the key to $L_A$ is given to $B$ and $C$, the key to $L_B$ is given to $A$ and $C$, and the key to $L_C$ is given to $A$ and $B$.

        Then each person is unable to open one of the locks, so no one can open the safe alone. However, any two people together can open all three locks. Therefore, the minimum is achieved by $\boxed{n=3,\ m=2}$.
    \end{solution}

    \begin{problem}{2022.3}
        Evaluate the integral
        \[
            \int_{-\infty}^{\infty}
            \frac{3}{(x-2)^2+7}\,dx.
        \]
    \end{problem}

    \begin{solution}
        Let
        \[
            u=x-2.
        \]
        Then the limits of integration remain $-\infty$ and $\infty$, so
        \[
            \int_{-\infty}^{\infty}
            \frac{3}{(x-2)^2+7}\,dx
            =
            3\int_{-\infty}^{\infty}
            \frac{1}{u^2+7}\,du.
        \]
        We use the standard formula
        \[
            \int_{-\infty}^{\infty}
            \frac{1}{u^2+a^2}\,du
            =
            \frac{\pi}{a}.
        \]
        Here $a=\sqrt{7}$, and therefore
        \[
            3\int_{-\infty}^{\infty}
            \frac{1}{u^2+7}\,du
            =
            3\cdot\frac{\pi}{\sqrt{7}}
            =
            \frac{3\pi}{\sqrt{7}}.
        \]
        Thus the answer is $\boxed{\dfrac{3\pi}{\sqrt{7}}}$.
    \end{solution}

    \begin{problem}{2022.4}
        A couple has two children.
        \begin{enumerate}
            \item[(a)] Given that one of the children is a daughter, what is the probability that the other child is also a daughter?
            \item[(b)] Given that one of the children is a daughter named Youngmi, what is the probability that the other child is also a daughter?
        \end{enumerate}
    \end{problem}

    \begin{solution}
        Let $G$ denote a girl and $B$ denote a boy. If we distinguish the two children by birth order, then the possible gender combinations are
        \[
            GG,\quad GB,\quad BG,\quad BB.
        \]

        For part (a), the condition that one child is a daughter means that at least one child is a girl. Thus the possible cases are
        \[
            GG,\quad GB,\quad BG.
        \]
        Among these three cases, only $GG$ has the property that the other child is also a daughter. Therefore,
        \[
            \Prob(\text{both are daughters}\mid \text{at least one is a daughter})
            =
            \frac{1}{3}.
        \]

        For part (b), the information that one child is a daughter named Youngmi is more specific. In the intended interpretation of the problem, this information identifies one particular child. Then the gender of the other child is still equally likely to be $G$ or $B$. Hence
        \[
            \Prob(\text{the other child is a daughter})
            =
            \frac{1}{2}.
        \]

        Therefore, the answers are $\boxed{\text{(a) } \dfrac{1}{3},\ \text{(b) } \dfrac{1}{2}}$.
    \end{solution}

    \begin{problem}{2022.5}
        Which of the following conjectures was not solved by Professor June Huh?
        \[
            \begin{array}{ll}
                (1) & \text{Read's conjecture} \\
                (2) & \text{Brylawski's conjecture} \\
                (3) & \text{Poincaré conjecture} \\
                (4) & \text{Rota's conjecture} \\
                (5) & \text{Okounkov's conjecture}
            \end{array}
        \]
    \end{problem}

    \begin{solution}
        Professor June Huh is known for connecting combinatorics with algebraic geometry and using these ideas to solve several major conjectures in combinatorics. His work is related to conjectures such as Read's conjecture, Brylawski's conjecture, Rota's conjecture, and Okounkov's conjecture.

        On the other hand, the Poincaré conjecture is a famous problem in $3$-dimensional topology. It was proved by Grigori Perelman using Ricci flow, not by June Huh.

        Therefore, the conjecture in the list that was not solved by June Huh is
        \[
            \boxed{(3)\ \text{Poincaré conjecture}}.
        \]
    \end{solution}

    \begin{problem}{2022.6}
        This mathematician, often called the father of linear programming, is famous for solving difficult problems while he was a doctoral student. One day, while studying at UC Berkeley, he arrived late to a class. The statistics professor Jerzy Neyman had written two famous unsolved problems on the blackboard, but the student thought they were homework problems. He took them home, solved them, and submitted complete solutions a few days later, thinking only that he had turned in the homework late. Who was this mathematician?
    \end{problem}

    \begin{solution}
        The mathematician described in the problem is George Dantzig.

        George Dantzig is known as one of the founders of linear programming and as the creator of the simplex method. For this reason, he is often called the father of linear programming.

        The anecdote in the problem is also a famous story about Dantzig. While he was a doctoral student at UC Berkeley, he arrived late to Jerzy Neyman's statistics class and copied down two problems written on the blackboard, believing that they were homework problems. He later submitted solutions, not realizing that the problems were actually open problems in statistics.

        Therefore, the answer is $\boxed{\text{George Dantzig}}$.
    \end{solution}

    \begin{problem}{2022.7}
        Two different coins are each tossed three times. Let $X$ and $Y$ be the numbers of heads obtained from the two coins, respectively. What is the probability that $X=Y$?
    \end{problem}

    \begin{solution}
        For each coin, the number of heads obtained in three tosses follows the binomial distribution
        \[
            X,Y\sim \Bin\left(3,\frac{1}{2}\right).
        \]
        Also, $X$ and $Y$ are independent. Therefore,
        \[
            \Prob(X=Y)
            =
            \sum_{k=0}^{3}\Prob(X=k)\Prob(Y=k).
        \]
        For each $k=0,1,2,3$, we have
        \[
            \Prob(X=k)
            =
            \binom{3}{k}\left(\frac{1}{2}\right)^3,
        \]
        and the same formula holds for $Y$. Hence
        \[
            \Prob(X=Y)
            =
            \sum_{k=0}^{3}
            \left(
                \binom{3}{k}\left(\frac{1}{2}\right)^3
            \right)^2
            =
            \frac{1}{64}
            \sum_{k=0}^{3}
            \binom{3}{k}^2.
        \]
        Thus
        \[
            \Prob(X=Y)
            =
            \frac{1}{64}
            \left(
                \binom{3}{0}^2+
                \binom{3}{1}^2+
                \binom{3}{2}^2+
                \binom{3}{3}^2
            \right)
            =
            \frac{1+9+9+1}{64}
            =
            \frac{20}{64}
            =
            \frac{5}{16}.
        \]
        Therefore, the answer is $\boxed{\dfrac{5}{16}}$.
    \end{solution}

    \begin{problem}{2023.1}
        In which year was the Fields Medal, won by Professor June Huh, first awarded?
    \end{problem}

    \begin{solution}
        The Fields Medal is one of the most prestigious awards in mathematics, and it is awarded by the International Mathematical Union. Professor June Huh received the Fields Medal in 2022.

        The Fields Medal was first awarded in 1936. The first two recipients were Lars Ahlfors and Jesse Douglas.

        Therefore, the answer is $\boxed{1936}$.
    \end{solution}

    \begin{problem}{2023.2}
        Consider the $6\times 6$ square grid with vertices $(0,0)$, $(6,0)$, $(0,6)$, and $(6,6)$. A path starts at $(0,0)$ and ends at $(6,6)$, moving only one unit to the right or one unit upward at each step. How many such paths never go above the line $x=y$? Equivalently, how many such paths stay in the region
        \[
            \{(x,y)\in\R^2:y\leq x\}?
        \]
    \end{problem}

    \begin{solution}
        To go from $(0,0)$ to $(6,6)$, a path must consist of $6$ right steps and $6$ upward steps. Without the condition $y\leq x$, the total number of such paths is
        \[
            \binom{12}{6}.
        \]

        We now count the paths that do go above the line $y=x$. By the reflection principle, the number of paths from $(0,0)$ to $(6,6)$ that at some point satisfy $y>x$ is
        \[
            \binom{12}{5}.
        \]
        Therefore, the number of paths that always stay in the region $y\leq x$ is
        \[
            \binom{12}{6}-\binom{12}{5}.
        \]
        Computing this gives
        \[
            \binom{12}{6}-\binom{12}{5}
            =
            924-792
            =
            132.
        \]
        Thus the answer is $\boxed{132}$.
    \end{solution}

    \begin{problem}{2023.3}
        A point on the surface of the Earth has the following property: starting from that point, one moves $5\mathrm{km}$ south, then $5\mathrm{km}$ east, and then $5\mathrm{km}$ north, and returns to the original point. How many such points are there on the surface of the Earth?

        Assume that the surface of the Earth is a perfect two-dimensional sphere, and that the North Pole and the South Pole are antipodal points. At the North Pole, the directions north, east, and west are not defined, and at the South Pole, the directions south, east, and west are not defined. Starting points for which one of these undefined directions occurs during the movement are excluded.
    \end{problem}

    \begin{solution}
        The North Pole is one such point. If we start at the North Pole, move $5\mathrm{km}$ south, then move $5\mathrm{km}$ east along a latitude circle, and finally move $5\mathrm{km}$ north, we return to the North Pole.

        There are also many such points near the South Pole. Let $Q$ be the point reached after moving $5\mathrm{km}$ south from the starting point. Suppose that the latitude circle passing through $Q$ has circumference
        \[
            \frac{5}{n}\mathrm{km}
        \]
        for some positive integer $n$. Then moving $5\mathrm{km}$ east from $Q$ goes exactly $n$ full times around that latitude circle and returns to $Q$. After that, moving $5\mathrm{km}$ north returns to the original starting point.

        Near the South Pole, the circumferences of latitude circles approach $0$. Hence, for every positive integer $n$, there is a latitude circle near the South Pole whose circumference is exactly $\dfrac{5}{n}\mathrm{km}$. For each such latitude circle, every point lying $5\mathrm{km}$ north of it gives a valid starting point.

        Therefore, there are not just one but infinitely many such points on the Earth's surface. In fact, this construction gives entire latitude circles of valid starting points, so there are uncountably many such points. For the usual quiz answer, the essential conclusion is
        \[
            \boxed{\text{infinitely many}}.
        \]
    \end{solution}

    \begin{problem}{2023.4}
        Let $f:(-1,1)\to\R$ be a continuous function. For each $\lambda>0$, define
        \[
            f_\lambda(x)
            :=
            \lambda f\left(\frac{x}{\lambda}\right),
        \]
        where $f_\lambda$ is defined for $x\in(-\lambda,\lambda)$.

        Suppose that for every sequence $\lambda_i\to\infty$, the functions $f_{\lambda_i}$ converge locally uniformly to $0$ on $\R$. Determine whether the following statement is true or false:
        \[
            f \text{ is differentiable at } 0 \text{ and } f'(0)=0.
        \]
    \end{problem}

    \begin{solution}
        We prove that the statement is true.

        First, evaluate at $x=0$. For every sequence $\lambda_i\to\infty$, we have
        \[
            f_{\lambda_i}(0)
            =
            \lambda_i f(0).
        \]
        Since $f_{\lambda_i}$ converges locally uniformly to $0$, in particular it converges pointwise to $0$ at $x=0$. Hence
        \[
            \lambda_i f(0)\to 0.
        \]
        This is possible for all sequences $\lambda_i\to\infty$ only if
        \[
            f(0)=0.
        \]

        Now we show that $f'(0)=0$. Let $h_i\to 0$ be any sequence with $h_i\neq 0$. For all sufficiently large $i$, we have $h_i\in(-1,1)$. Define
        \[
            \lambda_i=\frac{1}{|h_i|}.
        \]
        Then $\lambda_i\to\infty$, and
        \[
            h_i=\frac{\sgn(h_i)}{\lambda_i}.
        \]
        By the assumption, $f_{\lambda_i}\to 0$ locally uniformly on $\R$. In particular, this convergence is uniform on the compact set $\{-1,1\}$. Therefore,
        \[
            f_{\lambda_i}(\sgn(h_i))\to 0.
        \]
        But
        \[
            f_{\lambda_i}(\sgn(h_i))
            =
            \lambda_i
            f\left(
                \frac{\sgn(h_i)}{\lambda_i}
            \right)
            =
            \lambda_i f(h_i).
        \]
        Hence
        \[
            \lambda_i f(h_i)\to 0.
        \]
        Since $f(0)=0$, we get
        \[
            \left|
                \frac{f(h_i)-f(0)}{h_i}
            \right|
            =
            \left|
                \frac{f(h_i)}{h_i}
            \right|
            =
            \lambda_i |f(h_i)|
            \to 0.
        \]
        Thus
        \[
            \frac{f(h_i)-f(0)}{h_i}\to 0
        \]
        for every sequence $h_i\to 0$ with $h_i\neq 0$. Therefore, $f$ is differentiable at $0$ and
        \[
            f'(0)=0.
        \]

        Thus the answer is $\boxed{\text{True}}$.
    \end{solution}

    \begin{problem}{2024.1}
        When all natural numbers from $100$ to $999$ are written in a row, how many times does the digit $0$ appear?
    \end{problem}

    \begin{solution}
        Every natural number from $100$ to $999$ is a three-digit number.

        In a three-digit number, the hundreds digit cannot be $0$. Therefore, the digit $0$ can appear only in the tens digit or the units digit.

        First, count the numbers whose tens digit is $0$. The hundreds digit can be one of $1,2,\dots,9$, and the units digit can be one of $0,1,\dots,9$. Hence there are
        \[
            9\cdot 10=90
        \]
        such numbers.

        Next, count the numbers whose units digit is $0$. The hundreds digit can be one of $1,2,\dots,9$, and the tens digit can be one of $0,1,\dots,9$. Hence there are again
        \[
            9\cdot 10=90
        \]
        such numbers.

        Therefore, the total number of appearances of the digit $0$ is
        \[
            90+90=180.
        \]
        Thus the answer is $\boxed{180}$.
    \end{solution}

    \begin{problem}{2024.2}
        Which of the following people was born the latest?
        \[
            \begin{array}{ll}
                (1) & \text{John von Neumann} \\
                (2) & \text{Alan Turing} \\
                (3) & \text{John Nash} \\
                (4) & \text{Pierre de Fermat}
            \end{array}
        \]
    \end{problem}

    \begin{solution}
        We compare their years of birth:
        \[
            \begin{array}{ll}
                (1) & \text{John von Neumann: born in } 1903, \\
                (2) & \text{Alan Turing: born in } 1912, \\
                (3) & \text{John Nash: born in } 1928, \\
                (4) & \text{Pierre de Fermat: born in } 1607.
            \end{array}
        \]
        Among these four, the person born the latest is John Nash.

        Therefore, the answer is $\boxed{(3)\ \text{John Nash}}$.
    \end{solution}

    \begin{problem}{2024.4}
        Let $C$ be a convex subset of $\R^3$. Here convex means that for any two points $\mathbf{p},\mathbf{q}\in C$, the line segment joining $\mathbf{p}$ and $\mathbf{q}$ is also contained in $C$.

        Suppose that the following two facts are known:
        \begin{enumerate}
            \item[(i)] The intersection of $C$ with the $xy$-plane is
            \[
                C\cap\{z=0\}
                =
                \{(x,y,0):x^2+4y^2\leq 1\}.
            \]
            \item[(ii)] The set $C$ contains the $z$-axis.
        \end{enumerate}
        Can the volume of $C$ in the region $-1\leq z\leq 1$ be determined?
        \[
            \begin{array}{ll}
                (1) & \text{This cannot be determined.} \\
                (2) & \pi \\
                (3) & 2\pi \\
                (4) & 4\pi
            \end{array}
        \]
    \end{problem}

    \begin{solution}
        Let
        \[
            E=\{(x,y)\in\R^2:x^2+4y^2\leq 1\}.
        \]
        By condition (i),
        \[
            C\cap\{z=0\}=E\times\{0\}.
        \]
        The ellipse $E$ has semiaxes $1$ and $\dfrac{1}{2}$, so its area is
        \[
            \area(E)
            =
            \pi\cdot 1\cdot \frac{1}{2}
            =
            \frac{\pi}{2}.
        \]

        For each height $z$, define the horizontal cross-section
        \[
            C_z
            =
            \{(x,y)\in\R^2:(x,y,z)\in C\}.
        \]
        We claim that $C_z$ has the same area as $E$ for every $z$.

        First, we show that $C_z\subseteq E$. Suppose that $(x,y,z)\in C$ and
        \[
            x^2+4y^2>1.
        \]
        Since $C$ contains the $z$-axis, the point $(0,0,t)$ belongs to $C$ for every $t\in\R$. By convexity, the line segment joining $(x,y,z)$ and $(0,0,t)$ is contained in $C$.

        For any $\alpha\in(0,1)$, we may choose
        \[
            t=-\frac{\alpha z}{1-\alpha}
        \]
        so that this segment intersects the $xy$-plane at the point
        \[
            (\alpha x,\alpha y,0).
        \]
        Since $\alpha$ can be chosen sufficiently close to $1$, we may arrange that
        \[
            (\alpha x)^2+4(\alpha y)^2>1.
        \]
        This would give a point of $C\cap\{z=0\}$ lying outside $E\times\{0\}$, contradicting condition (i). Therefore,
        \[
            C_z\subseteq E
        \]
        for every $z$.

        Conversely, we show that the interior of $E$ is contained in every $C_z$. Take $(x,y)$ with
        \[
            x^2+4y^2<1.
        \]
        Choose $\alpha\in(0,1)$ sufficiently close to $1$ so that
        \[
            \left(\frac{x}{\alpha}\right)^2
            +
            4\left(\frac{y}{\alpha}\right)^2
            \leq 1.
        \]
        Then
        \[
            \left(\frac{x}{\alpha},\frac{y}{\alpha},0\right)\in C.
        \]
        Also, since $C$ contains the $z$-axis,
        \[
            \left(0,0,\frac{z}{1-\alpha}\right)\in C.
        \]
        By convexity, their convex combination also belongs to $C$, and
        \[
            \alpha
            \left(\frac{x}{\alpha},\frac{y}{\alpha},0\right)
            +
            (1-\alpha)
            \left(0,0,\frac{z}{1-\alpha}\right)
            =
            (x,y,z).
        \]
        Hence $(x,y)\in C_z$.

        Thus, for every $z$,
        \[
            \operatorname{int}(E)\subseteq C_z\subseteq E.
        \]
        Since the boundary of the ellipse has area $0$, we have
        \[
            \area(C_z)=\area(E)=\frac{\pi}{2}
        \]
        for every $z$. Therefore, the volume of $C$ in the region $-1\leq z\leq 1$ is
        \[
            \int_{-1}^{1}\area(C_z)\dd z
            =
            \int_{-1}^{1}\frac{\pi}{2}\dd z
            =
            \pi.
        \]

        Therefore, the answer is $\boxed{(2)\ \pi}$.
    \end{solution}

    \begin{problem}{2025.1}
        A fair coin is tossed repeatedly until either two heads in a row or two tails in a row appear. On average, how many tosses are needed before the process stops?
    \end{problem}

    \begin{solution}
        On the first toss, it is impossible to have two consecutive equal outcomes. After the first toss, the process stops exactly when the next toss gives the same result as the previous toss.

        Since the coin is fair, from the second toss onward we have
        \[
            \Prob(\text{the new toss matches the previous toss})
            =
            \frac{1}{2}.
        \]
        Therefore, after the first toss, the number of additional tosses needed until the process stops follows a geometric distribution with success probability $\dfrac{1}{2}$. Its expected value is
        \[
            \frac{1}{1/2}=2.
        \]
        Since the first toss is always made, the expected total number of tosses is
        \[
            1+2=3.
        \]
        Thus the answer is $\boxed{3}$.
    \end{solution}

    \begin{problem}{2025.2}
        The nineteenth-century mathematicians Cauchy and Weierstrass established rigorous definitions of limits. Who was the first mathematician to prove that the sum of the reciprocals of the squares of all positive integers is
        \[
            \sum_{n=1}^{\infty}\frac{1}{n^2}
            =
            \frac{\pi^2}{6}?
        \]
        \[
            \begin{array}{ll}
                (1) & \text{David Hilbert} \\
                (2) & \text{Georg Cantor} \\
                (3) & \text{Bernhard Riemann} \\
                (4) & \text{Leonhard Euler}
            \end{array}
        \]
    \end{problem}

    \begin{solution}
        The problem of finding the exact value of the infinite series
        \[
            \sum_{n=1}^{\infty}\frac{1}{n^2}
        \]
        is known as the Basel problem.

        This problem had been known since the seventeenth century, and it was Leonhard Euler who first found and proved the famous identity
        \[
            \sum_{n=1}^{\infty}\frac{1}{n^2}
            =
            \frac{\pi^2}{6}.
        \]
        Therefore, the answer is $\boxed{(4)\ \text{Leonhard Euler}}$.
    \end{solution}

    \begin{problem}{2025.3}
        Find the equation of the smallest sphere tangent to all four planes
        \[
            x=0,\qquad y=0,\qquad z=0,\qquad x+y+z=3.
        \]
    \end{problem}

    \begin{solution}
        Since the sphere is tangent to the three coordinate planes
        \[
            x=0,\qquad y=0,\qquad z=0,
        \]
        and we are looking for the smallest such sphere, its center must lie in the first octant and have the form
        \[
            \mathbf{c}=(a,a,a)
        \]
        by symmetry. The distance from $\mathbf{c}$ to each coordinate plane is $a$, so the radius of the sphere is
        \[
            r=a.
        \]

        The sphere must also be tangent to the plane
        \[
            x+y+z=3.
        \]
        The distance from $(a,a,a)$ to this plane is
        \[
            \frac{|a+a+a-3|}{\sqrt{1^2+1^2+1^2}}
            =
            \frac{3-3a}{\sqrt{3}},
        \]
        since the smallest sphere lies between the coordinate planes and the plane $x+y+z=3$. This distance must equal the radius $a$, so
        \[
            \frac{3-3a}{\sqrt{3}}=a.
        \]
        Hence
        \[
            3-3a=\sqrt{3}a,
        \]
        and therefore
        \[
            a
            =
            \frac{3}{3+\sqrt{3}}
            =
            \frac{3-\sqrt{3}}{2}.
        \]

        Thus the center is
        \[
            \left(
                \frac{3-\sqrt{3}}{2},
                \frac{3-\sqrt{3}}{2},
                \frac{3-\sqrt{3}}{2}
            \right),
        \]
        and the radius is
        \[
            \frac{3-\sqrt{3}}{2}.
        \]
        Therefore, the equation of the sphere is
        \[
            \boxed{
            \left(x-\frac{3-\sqrt{3}}{2}\right)^2
            +
            \left(y-\frac{3-\sqrt{3}}{2}\right)^2
            +
            \left(z-\frac{3-\sqrt{3}}{2}\right)^2
            =
            \left(\frac{3-\sqrt{3}}{2}\right)^2
            }.
        \]
    \end{solution}

    \part{Expected Problems}

\end{document}